% ============================================================
%  Дипломски рад — Кристина Живковић
%  Факултет организационих наука
%
%  КОМПАЈЛИРАЊЕ (обавезно XeLaTeX, не pdfLaTeX):
%      latexmk -xelatex rad.tex
%  или, ако latexmk није при руци:
%      xelatex rad   →   biber rad   →   xelatex rad   →   xelatex rad
% ============================================================

\documentclass[12pt,a4paper,oneside]{report}

% ============================================================
%  Преамбула — заједничка подешавања за цео рад
%  Компајлира се искључиво са XeLaTeX (због fontspec/unicode-math)
% ============================================================

% ── језик и фонтови ─────────────────────────────────────────
\usepackage{fontspec}
\usepackage{polyglossia}
\setmainlanguage{serbian}
\setmainfont{Times New Roman}

% Моноспејс фонт мора да садржи ћирилицу (подразумевани Latin Modern Mono
% је нема, што polyglossia пријављује као грешку код \texttt{...}).
\setmonofont{Consolas}[Scale=0.92]

% ── математика ──────────────────────────────────────────────
% ВАЖНО: amsmath иде ПРЕ unicode-math, а amssymb се НЕ учитава --
% unicode-math већ дефинише све те симболе, па би amssymb изазвао
% сукоб („Command \eth already defined“).
\usepackage{amsmath}
\usepackage{unicode-math}
\setmathfont{Cambria Math}
\allowdisplaybreaks          % дозволи прелом вишередних формула преко стране

% ── распоред стране ─────────────────────────────────────────
% Маргине по шаблону ФОН-а: горња, доња и десна 2,5 цм; лева 3,1 цм.
\usepackage[a4paper,left=3.1cm,right=2.5cm,top=2.5cm,bottom=2.5cm]{geometry}
\usepackage{setspace}
\onehalfspacing               % проред између редова 1,5

% Проред параграфа по шаблону ФОН-а: 12pt пре и 6pt након параграфа.
% Word сажима суседне размаке на ВЕЋИ од та два (12pt), не на њихов
% збир, па је 12pt тачан еквивалент. LaTeX има само један размак
% између параграфа (\parskip), па се он поставља на 12pt.
% Растезање (plus/minus) намерно је нула: без тога LaTeX додатно
% развлачи размаке да попуни страну, што ствара неуједначене празнине.
\setlength{\parindent}{0pt}
\setlength{\parskip}{12pt}

% Обострано поравнање је подразумевано у LaTeX-у; овде се само
% ублажава превише агресивно растезање речи.
\usepackage{microtype}

% Латиничне речи не смеју да се преламају (види rad.tex), па TeX у понеком
% реду нема ниједно место прелома. \emergencystretch му дозвољава да такав
% ред мало растегне уместо да га пусти преко десне маргине.
\setlength{\emergencystretch}{1em}

% Српски обрасци преламања примењују се и на латиничне речи, а не умеју да
% их преломе на правом месту (discove-ry, le-ast, AR-FIMA, XG-Boost). У
% ћириличном тексту је свака латинична реч страна, па се преламање свих
% таквих речи забрањује — реч наведена у \hyphenation без иједне цртице
% нема ниједно дозвољено место прелома. Списак је извучен из тела рада.
% Списак стоји у rad.tex одмах иза \begin{document}, а не овде: polyglossia
% активира српски тек по уласку у документ, па изузетак уписан у преамбули
% (па и кроз \AtBeginDocument) припадне другом језику и остане без дејства.

% ── уједначавање вертикалних размака ────────────────────────
% Пошто је \parskip постављен на 12pt, LaTeX-ови подразумевани
% размаци око набрајања, формула и наслова СЕ ДОДАЈУ на њега и
% стварају видљиве празне редове. Зато се сви експлицитно своде
% на исти размак од 12pt, тако да документ има један једини
% вертикални ритам.

% Набрајања (itemize, enumerate, description)
\usepackage{enumitem}
\setlist{
  topsep    = 0pt,     % изнад/испод целе листе — наслеђује \parskip
  partopsep = 0pt,
  parsep    = 0pt,
  itemsep   = 6pt,     % између ставки: пола размака параграфа
  listparindent = 0pt
}

% Формуле: amsmath подразумевано додаје сопствени размак изнад и
% испод, који се сабира са \parskip. Своди се на исту вредност.
% Поставља се кроз \AtBeginDocument јер amsmath своје вредности
% дефинише при учитавању величине фонта и прегазио би ранију поставку.
% Формула је наставак реченице која јој претходи, а не нов параграф,
% па добија 6pt — пола размака параграфа. Са 12pt би изгледала као
% да је одвојена празним редом, што је управо оно што се избегава.
\AtBeginDocument{%
  \setlength{\abovedisplayskip}{6pt}%
  \setlength{\belowdisplayskip}{6pt}%
  \setlength{\abovedisplayshortskip}{6pt}%
  \setlength{\belowdisplayshortskip}{6pt}%
}

% Плутајући објекти (табеле): размак изнад и испод натписа и
% између објекта и текста.
\setlength{\textfloatsep}{12pt}
\setlength{\floatsep}{12pt}
\setlength{\intextsep}{12pt}
\setlength{\abovecaptionskip}{6pt}
\setlength{\belowcaptionskip}{6pt}

% ── табеле ──────────────────────────────────────────────────
\usepackage{booktabs}
\usepackage{multirow}
\usepackage{needspace}   % држи пасус и његову једначину на истој страни
\usepackage{xurl}        % дозвољава прелом дугих идентификатора у \url
\usepackage{tabularx}
\usepackage{longtable}
\usepackage{array}
\newcolumntype{L}[1]{>{\raggedright\arraybackslash}p{#1}}

% ── наслови по шаблону ФОН-а ────────────────────────────────
% Ниво 1: СВИМ ВЕЛИКИМ СЛОВИМА, фонт 14, затамњено.
% Ниво 2: СВИМ ВЕЛИКИМ СЛОВИМА, фонт 14 (без задебљања).
% Ниво 3: малим словима, фонт 14.
% Поглавља се не преламају на нову страну (шаблон нема преломе),
% па се користи \chapter* логика преко titlesec-а без \clearpage.
\usepackage{titlesec}

% Размак ИСПОД наслова постављен је на 0pt, а не на 6pt: \parskip од
% 12pt већ увлачи размак пре првог параграфа иза наслова, па би 6pt
% на то дало 18pt и видљиву празнину. Размак ИЗНАД наслова је 6pt,
% што уз \parskip претходног параграфа даје укупно 18pt раздвајања
% између претходног текста и новог наслова.
\titleformat{\chapter}[hang]
  {\fontsize{14}{17}\selectfont\bfseries\MakeUppercase}
  {\thechapter}{1em}{}
\titlespacing*{\chapter}{0pt}{0pt}{0pt}

\titleformat{\section}[hang]
  {\fontsize{14}{17}\selectfont\MakeUppercase}
  {\thesection}{1em}{}
\titlespacing*{\section}{0pt}{6pt}{0pt}

\titleformat{\subsection}[hang]
  {\fontsize{14}{17}\selectfont}
  {\thesubsection}{1em}{}
\titlespacing*{\subsection}{0pt}{6pt}{0pt}

% \paragraph се користи за истакнуте поднаслове у тексту (нпр. „Х1 --- ...“).
% Мора остати у истом реду са текстом који следи, без додатног размака.
\titleformat{\paragraph}[runin]
  {\normalfont\bfseries}{}{0pt}{}
\titlespacing*{\paragraph}{0pt}{0pt}{1em}

% Поглавље почиње на новој страни, али без празне стране између
% (oneside документ то и иначе поштује).
\let\originalchapter\chapter

% ── референце и хипервезе ───────────────────────────────────
% ── размаци у садржају ──────────────────────────────────────
% Опција [titles] оставља обликовање наслова класи, а даје контролу само
% над вертикалним размацима између ставки садржаја. Без ње би tocloft
% преузео и изглед наслова поглавља, што би поништило подешавања titlesec-а.
\usepackage[titles]{tocloft}
\setlength{\cftbeforechapskip}{6pt}      % мали додатни размак пред поглављима
\setlength{\cftbeforesecskip}{4pt}       % између ставки нивоа x.y
\setlength{\cftbeforesubsecskip}{0pt}    % између ставки нивоа x.y.z
\setlength{\cftbeforesubsubsecskip}{0pt}

% ── размаци око истакнутих формула ──────────────────────────
% Изједначени размак изнад и испод приказа, мањи од подразумеваног.
% Иде кроз \AtBeginDocument јер setspace при \onehalfspacing поставља
% ове дужине на почетку документа, па би ранија вредност била прегажена.
\AtBeginDocument{%
  \setlength{\abovedisplayskip}{6pt plus 2pt minus 2pt}%
  \setlength{\belowdisplayskip}{6pt plus 2pt minus 2pt}%
  \setlength{\abovedisplayshortskip}{3pt plus 1pt}%
  \setlength{\belowdisplayshortskip}{3pt plus 1pt}}

\usepackage[hidelinks]{hyperref}
\usepackage{caption}

% Натписи по шаблону ФОН-а: „Табела 1: ...“ изнад табеле,
% „Слика 1. ...“ испод слике. Фонт исти као у остатку рада (12).
\captionsetup{font=normalsize,labelfont=bf,labelsep=period}
\captionsetup[table]{labelsep=colon,justification=raggedright,singlelinecheck=false}
\captionsetup[figure]{justification=centering}
\addto\captionsserbian{%
  \renewcommand{\tablename}{Табела}%
  \renewcommand{\figurename}{Слика}%
}

% ── непрекидна нумерација (1), (2), (3)... по шаблону ───────
% Шаблон приказује формуле нумерисане просто бројем у загради,
% а табеле и слике редним бројем кроз цео рад — не по поглављу.
% Зато се \numberwithin НЕ користи: подразумевана нумерација у
% класи report јесте по поглављу, па се експлицитно враћа на
% непрекидну.
\counterwithout{equation}{chapter}
\counterwithout{table}{chapter}
\counterwithout{figure}{chapter}

% ── скраћенице које се често понављају ──────────────────────
\newcommand{\RV}{\mathrm{RV}}
\newcommand{\RQ}{\mathrm{RQ}}
\newcommand{\BPV}{\mathrm{BPV}}
\newcommand{\MedRV}{\mathrm{MedRV}}
\newcommand{\TrBPV}{\mathrm{TrBPV}}
\newcommand{\TPQ}{\mathrm{TPQ}}
\newcommand{\RS}{\mathrm{RS}}
\newcommand{\SJ}{\mathrm{SJ}}
\newcommand{\QV}{\mathrm{QV}}
\newcommand{\QLIKE}{\mathrm{QLIKE}}
\newcommand{\RMSE}{\mathrm{RMSE}}
\newcommand{\MAE}{\mathrm{MAE}}
\newcommand{\VaR}{\mathrm{VaR}}
\newcommand{\DM}{\mathrm{DM}}
\newcommand{\BIC}{\mathrm{BIC}}
\newcommand{\VIX}{\mathrm{VIX}}
\newcommand{\Hit}{\mathrm{Hit}}
\newcommand{\med}{\operatorname{med}}
% општа ознака естиматора отпорног на скокове (један ниво фактора А2)
\newcommand{\IVhat}{\widehat{\mathrm{IV}}}

% ── слике и ознаке за вредности из новог покретања ──────────
\usepackage{graphicx}
\usepackage{xcolor}
\usepackage{tikz}
\usetikzlibrary{arrows.meta,positioning,fit,backgrounds}

% Стилови за дијаграм истраживачког тока (слика 1).
\tikzset{
  faza/.style     = {draw, rounded corners=2pt, align=center, inner sep=6pt,
                     font=\small, text width=3.05cm, minimum height=1.25cm,
                     fill=black!4},
  detalj/.style   = {align=center, font=\footnotesize\itshape, text width=3.05cm},
  tok/.style      = {-{Stealth[length=2.6mm]}, thick},
  medjukorak/.style = {align=center, font=\footnotesize, text width=2.1cm}
}
\newcommand{\todo}[1]{\textcolor{red}{[\textbf{TODO:} #1]}}
% оквир на месту слике која ће бити генерисана после покретања
\newcommand{\slikaplaceholder}[1]{%
  \fbox{\parbox[c][5cm][c]{0.85\textwidth}{\centering\todo{#1}}}}


% ── библиографија: АПА стил (захтев шаблона ФОН-а) ──────────
% Литература се ређе азбучно по презимену првог аутора.
% maxcitenames=3: у цитату се наводе сва три аутора (Bollerslev, Patton
% & Quaedvlieg, 2016) уместо APA скраћивања на „Bollerslev и др.“.
\usepackage[backend=biber,style=apa,sorting=nyt,maxbibnames=99,
            maxcitenames=3,mincitenames=1,
            language=english,autolang=none]{biblatex}
% biblatex-apa нема српску локализацију. Пошто је главни језик рада
% српски, biblatex учитава `serbian.lbx`, у коме APA датумски макрои
% не постоје, па компајлирање пуца на „Undefined control sequence“.
% Решење: српски се мапира на енглеску APA датотеку, а недостајући
% макрои се дефинишу преузимањем из `english-apa.lbx`.
\DeclareLanguageMapping{serbian}{english-apa}
\DeclareLanguageMapping{serbianc}{english-apa}
\DeclareLanguageMapping{english}{english-apa}
\addbibresource{literatura.bib}

% Списак литературе се слаже без обостраног поравнања: при поравнању LaTeX
% растеже размаке да попуни ред, што дуге DOI адресе разбија на „https : / /“.
\AtBeginBibliography{\raggedright}

% Компатибилносни слој: APA датумски макрои које српска датотека нема.
\makeatletter
\providecommand{\mkbibdateapalong}[3]{%
  \iffieldundef{#3}{}{\stripzeros{\thefield{#3}}\nobreakspace}%
  \iffieldundef{#2}{}{\mkbibmonth{\thefield{#2}}\iffieldundef{#3}{}{}}%
  \iffieldundef{#1}{}{\iffieldundef{#2}{}{, }\mkbibdatelong@year{#1}}}
\providecommand{\mkbibdatelong@year}[1]{\thefield{#1}}
\providecommand{\mkbibdateapalongextra}[4]{%
  \iffieldundef{#1}{}{\thefield{#1}}%
  \iffieldundef{#3}{}{, \mkbibmonth{\thefield{#3}}}%
  \iffieldundef{#4}{}{\nobreakspace\stripzeros{\thefield{#4}}}}
\providecommand{\mkbibdateapashort}[3]{%
  \iffieldundef{#1}{}{\thefield{#1}}}
\makeatother

% Шаблон ФОН-а тражи српске везнике у цитатима:
%   два аутора      → (Влајић и Лечић-Цветковић, 2021)
%   више од два     → (Јовић и др, 2021)
% Стрингови се дефинишу за ENGLISH, јер је то језик библиографије.
\NewBibliographyString{jourarticle}
\DefineBibliographyStrings{english}{%
  andothers = {и~др\adddot},
  and       = {и},
  in        = {In},
  page      = {p\adddot},
  pages     = {pp\adddot},
  % APA испис броја чланка („Article 14“) — српска локализација нема овај низ
  jourarticle = {Article},
}
% Исти низови морају бити дефинисани и за српски, јер biblatex бира низове
% према активном језику документа (serbianc), а не према језику библиографије.
\DefineBibliographyStrings{serbian}{%
  andothers = {и~др\adddot},
  and       = {и},
  in        = {In},
  jourarticle = {Article},
}
\DefineBibliographyStrings{serbianc}{%
  andothers = {и~др\adddot},
  and       = {и},
  in        = {In},
  jourarticle = {Article},
}
% APA стил користи амперсанд (&) и у заградним цитатима и у списку
% литературе; шаблон ФОН-а тражи српско „и“, па се замењује свуда.
% Новији biblatex делимитере поставља преко \DeclareDelimFormat, па их
% \renewcommand више не мења: APA контекст би вратио амперсанд.
% Декларације иду кроз \AtBeginDocument јер biblatex сам стил (apa.bbx, apa.cbx)
% учитава тек на крају преамбуле, па би ранија дефиниција била прегажена.
% APA пред последњим аутором ставља зарез („A, B, & C“); уз српско „и“ зарез се
% изоставља, па се Оксфордски зарез гаси.
\AtBeginDocument{%
  \DeclareDelimFormat{finalnamedelim}{\addspace\bibstring{and}\addspace}%
  \DeclareDelimFormat[parencite,citation]{finalnamedelim}{%
    \addspace\bibstring{and}\addspace}%
  \DeclareDelimFormat[bib,biblist]{finalnamedelim}{%
    \addspace\bibstring{and}\addspace}%
  \DeclareDelimFormat[textcite]{finalnamedelim}{%
    \addspace\bibstring{and}\addspace}%
  \DeclareDelimFormat{andothersdelim}{\addspace}%
  \renewcommand*{\finalandcomma}{}%
}
% Српски везник „и“ не тражи зарез испред себе, док APA пред „&“ ставља зарез.
% Датотека језика (english-apa.lbx) враћа \finalandcomma при свакој промени
% језика, па се гаси непосредно пре исписа цитата и списка литературе.
% Зарез се враћа при свакој промени језика (свака одредница у списку литератури
% носи језик), па се гаси у самим језичким додацима, а не једном у преамбули.
\DefineBibliographyExtras{english}{\renewcommand*{\finalandcomma}{}}
\DefineBibliographyExtras{serbian}{\renewcommand*{\finalandcomma}{}}
\DefineBibliographyExtras{serbianc}{\renewcommand*{\finalandcomma}{}}

% Зарез пред последњим аутором у списку литературе НЕ долази од делимитера него од
% самог формата имена: APA после сваког неконачног обрнутог имена („Презиме, И.“)
% додаје зарез (apa.bbx, \apablx@ifrevnameappcomma), па се уз српско „и“ добија
% „Барндорф-Нилсен, О. Е., и Шепард, Н.“. Зато се услов помера једно име раније:
% зарез се изоставља и на претпоследњем имену, дакле непосредно пред „и“, док
% зарези између осталих имена остају. Цитати се не мењају, јер у њима имена нису
% обрнута, па се овај зарез тамо и не додаје.
\makeatletter
\AtBeginDocument{%
  \renewcommand*{\apablx@ifrevnameappcomma}{%
    \iftoggle{apablx@wantcomma}
      {\ifboolexpr{
         test {\ifnumless{\value{listcount}}{\the\numexpr\value{liststop}-1\relax}}
         or
         test \ifmorenames
       }}
      {\@secondoftwo}}%
}
\makeatother

\title{Дизајн мерења наспрам избора модела у предвиђању реализоване волатилности}
\author{Кристина Живковић}
\date{2026.}

\begin{document}

% ── насловна страна (по шаблону ФОН-а) ──────────────────────
\begin{titlepage}
\centering
\vspace*{1.5cm}

{\bfseries УНИВЕРЗИТЕТ У БЕОГРАДУ\par}
\vspace{0.4em}
{\bfseries ФАКУЛТЕТ ОРГАНИЗАЦИОНИХ НАУКА\par}

\vspace{5cm}

{\bfseries ЗАВРШНИ РАД\par}

\vspace{1cm}

{\fontsize{18}{22}\selectfont\bfseries
Утицај дизајна мерења и избора модела\\[0.3em]
на тачност предвиђања реализоване\\[0.3em]
волатилности\par}

\vspace{4cm}

\begin{minipage}[t]{0.48\textwidth}
\raggedright
{\bfseries Ментор}\\[0.8em]
{\bfseries Др Александар Ракићевић,}\\
{\bfseries Доцент}
\end{minipage}
\hfill
\begin{minipage}[t]{0.48\textwidth}
\raggedright
{\bfseries Студент}\\[0.8em]
{\bfseries Кристина Живковић}\\
{\bfseries 2020/0029}
\end{minipage}

\vfill

{\bfseries Београд, 2026. године\par}
\end{titlepage}

% Насловна страна, апстракт и садржај се не нумеришу; нумерација
% почиње од прве стране увода (захтев шаблона).
\pagenumbering{gobble}

% ── апстракт ────────────────────────────────────────────────
\chapter*{АПСТРАКТ}
\addcontentsline{toc}{chapter}{Апстракт}

Реализована волатилност представља једну од централних величина савременог управљања
финансијским ризиком, будући да се на њој заснивају мере попут вредности при ризику,
а користи се и при вредновању опција. За разлику од приноса, чија је предвидљивост емпиријски занемарљива,
интензитет кретања цена показује изражену перзистентност и стога представља проблем у
коме стваран сигнал постоји. Литература о предвиђању реализоване волатилности развија
се у два тока: један пореди моделе при фиксираним мерама, од HAR породице преко
стабала одлучивања до архитектура дубоког учења, а други пореди саме реализоване мере
при фиксираном моделу. Реализована волатилност је, међутим, изведена, а не посматрана
величина, чија свака вредност произлази из низа методолошких одлука.

Рад испитује да ли разлике у тачности предвиђања реализоване волатилности потичу
претежно од избора модела или од одлука донетих при мерењу улазних величина.

Спроведен је факторски експеримент у коме се дизајн мерења мења по три
димензије: фреквенцији узорковања, естиматору отпорном на скокове и тесту за детекцију
скокова. Свака варијанта дизајна мерења комбинује се са скупом од шеснаест
економетријских модела, модела машинског учења и претходно обучених модела опште
намене за временске серије. Модели се поново оцењују сваких 21 дан трговања на
проширујућем прозору који почиње 2010. године, а вреднују се у периоду 2015–2025.
помоћу QLIKE функције губитка, Диболд–Маријановог теста, скупа поверења модела и
тестова вредности при ризику. Удео варијансе губитка који објашњава сваки фактор
утврђује се двема декомпозицијама варијансе, са блокирањем по симболу и режиму и
кластер-робусним стандардним грешкама: једном за моделе чији улази зависе од свих
димензија мерења и другом за моделе који зависе само од фреквенције узорковања.

Резултати показују да дизајн мерења објашњава удео варијансе губитка који је мањи од
удела избора модела, али истог реда величине, и то на једнодневном хоризонту, док међу
моделима који као улаз користе само реализовану варијансу удео мерења удео модела и
надмашује; готово цео тај ефекат носи фреквенција узорковања, а не естиматор или тест
за детекцију скокова. Информација о скоковима носи стваран али мали сигнал: разлика у
односу на референтни модел није значајна на референтном дизајну, али јесте на већини
дизајна при једнодневном хоризонту, а плацебо провера показује да потиче од стварних
дана са скоком. Рангирање модела осетно се мења при преласку на економску евалуацију,
уз Кендалов коефицијент $\tau$ између 0,300 и 0,400. Налази подупиру тврдњу да дизајн
мерења треба наводити и образлагати као део методологије, једнако као архитектуру
модела.

\vspace{1em}
\noindent\textbf{Кључне речи:} реализована волатилност, дизајн мерења, HAR модел,
детекција скокова, декомпозиција варијансе, вредност при ризику.

% ── садржај ─────────────────────────────────────────────────
\renewcommand{\contentsname}{САДРЖАЈ}
\tableofcontents

% ── поглавља ────────────────────────────────────────────────
% Нумерација страница почиње од прве стране увода (шаблон ФОН-а).
\clearpage
\pagenumbering{arabic}
\setcounter{page}{1}

\chapter{Увод}

Развој метода машинског и дубоког учења омогућио је анализу сложених појава чија
структура измиче традиционалним статистичким приступима. У области финансија,
међутим, примена ових метода наилази на границу која није алгоритамске, него
структурне природе. Гу, Кели и Сју \parencite{guKellyXiu2020} упоредили су читав
репертоар метода машинског учења на скупу од преко деветсто предиктора и утврдили да
и најбоље међу њима објашњавају мање од пола процента варијансе месечних приноса.
Тај налаз указује на то да правац кретања цена није проблем са скривеним решењем,
већ проблем у коме сигнала готово да нема.

Постоји, међутим, изразита асиметрија у предвидљивости финансијских временских
серија. Док је аутокорелација приноса практично једнака нули, аутокорелација
квадрираних приноса снажно је позитивна и споро опада, остајући значајна и након више
десетина дана трговања. Другим речима, не зна се да ли ће цена расти или падати,
али се са знатном тачношћу може предвидети колико ће се снажно кретати. Ова асиметрија
има економско објашњење: предвидљивост правца сама себе уништава, јер је арбитражери
искоришћавају чим се појави, док предвидљивост интензитета не носи директну опкладу
којом би била избрисана. Реализована волатилност, мерена сумом квадрата интрадневних
логаритамских приноса, стога представља ретку величину у финансијама која има стваран
и постојан сигнал, а истовремено и величину на којој почива управљање ризиком: из ње
се изводи вредност при ризику, која је део регулаторних захтева за капитал, а користи
се и при вредновању опција.

Модели за предвиђање реализоване волатилности развијали су се у две паралелне линије.
Економетријску линију утемељио је Корси \parencite{corsi2009} хетерогеним
ауторегресивним моделом, који дуготрајно памћење волатилности апроксимира каскадом
компоненти на дневном, недељном и месечном хоризонту, инспирисаном хипотезом о
хетерогеном тржишту \parencite{muller1997}, дакле идејом да учесници реагују различитим
брзинама. Проширења овог модела уводе разлагање на континуирану и скоковиту компоненту
\parencite{abd2007}, док Патон и Шепард \parencite{pattonSheppard2015} показују да је
негативна семиваријанса знатно значајнија за будућу волатилност од позитивне, и да
негативни скокови воде вишој, а позитивни нижој будућој волатилности. Друга линија
примењује методе машинског учења: Кристенсен, Сигор и Велијев
\parencite{christensen2023} налазе да алгоритми машинског учења надмашују HAR
породицу чак и када располажу само истим улазима као HAR модел, дакле дневним,
недељним и месечним просецима прошле волатилности, при чему су добици највећи на
дужим хоризонтима.

Актуелна расправа концентрисана је на архитектуре дубоког учења, а свеобухватан преглед
области \parencite{leushuisPetkov2026}, који обухвата поређење модела у литератури од
2000. до прве половине 2024. године, показује да питање остаје неразрешено: налази се
разилазе и не дају јединствен одговор о томе да ли те архитектуре надмашују
једноставније моделе.

У већини ових радова реализоване мере узимају се као дате, а варира се модел.
Реализована волатилност, међутим, није посматрана него изведена величина: свака њена
вредност резултат је низа методолошких одлука, дакле избора
фреквенције узорковања на којој се рачунају приноси, избора естиматора отпорног на
скокове којим се процењује континуирани део волатилности, и избора статистичког теста
којим се дан проглашава даном са скоком. Ниједна од тих одлука није неутрална: превисока
фреквенција узорковања уводи микроструктурни шум који систематски надувава меру, а
различити тестови за детекцију скокова дају битно различите скупове детектованих
догађаја. Лиу, Патон и Шепард \parencite{lps2015} показали су да избор реализоване
мере мења њену тачност као проксија, али су при томе држали модел фиксним. Да грешка
мерења заиста утиче на прогнозу, показали су Болерслев, Патон и Кведвлиг
\parencite{bpq2016} HARQ моделом, који коефицијент уз дневни регресор интерагује са
реализованом квартичношћу управо да би узео у обзир временски променљиву грешку
мерења. Тај модел, међутим, грешку мерења моделује унутар једног модела, уместо да је
третира као димензију која се може експериментално варирати.

Овај рад зато мерење и модел третира као две димензије истог експеримента: обе се
варирају на истом панелу америчких акција, у потпуном факторском плану, а њихови
доприноси тачности раздвајају се декомпозицијом варијансе.

\section{Циљ рада}

Општи циљ рада јесте да се измери колики удео разлика у тачности предвиђања реализоване
волатилности потиче од избора модела, а колики од одлука донетих при мерењу улазних
величина.

Конкретни циљеви рада су:

\begin{itemize}
\item имплементирати сопствени систем за конструкцију реализованих мера, у коме се
      дизајн мерења може мењати по димензијама, без измене осталих делова поступка;
\item спровести факторски експеримент у коме се тридесет шест дизајна мерења укршта са
      шеснаест модела, кроз три хоризонта предвиђања, уз поновно оцењивање модела
      сваких 21 дан трговања на проширујућем прозору;
\item измерити тачност предвиђања функцијом губитка робусном на шум у проксију, уз
      тестове статистичке значајности и скуп поверења модела;
\item бројчано раздвојити удео варијансе губитка који се приписује дизајну мерења од
      удела који се приписује избору модела, помоћу две декомпозиције варијансе: једне
      за моделе чији улази зависе од свих трију димензија мерења и друге за моделе
      чији улази зависе само од фреквенције узорковања;
\item утврдити да ли информација о скоковима доприноси тачности, како тај допринос
      зависи од хоризонта предвиђања и да ли опстаје када се индикатори скокова
      насумично пермутују;
\item утврдити да ли се рангирање модела мења при преласку са статистичке на економску
      евалуацију, путем тестова вредности при ризику.
\end{itemize}

\section{Методологија}

За потребе рада развијен је сопствени систем за конструкцију реализованих мера и
оцењивање модела, у програмском језику Python. Тачност се мери QLIKE функцијом губитка,
изабраном зато што је Патон \parencite{patton2011} показао да већина уобичајених
функција губитка даје погрешно рангирање модела када је прокси за стварну волатилност
зашумљен, а да је QLIKE међу малобројним функцијама робусним на ту појаву. Значајност
разлика утврђује се упареним тестом предиктивне тачности и скупом поверења модела, уз
контролу стопе лажних открића, а доприноси појединачних фактора раздвајају се
декомпозицијом варијансе. Мере, правила искључивања и план статистичке анализе
записани су пре коначног покретања експеримента, а параметри који нису предмет
варирања фиксирани су унапред и наведени уз образложење; одступања од тог редоследа
наведена су на месту на коме се појављују. Поступак је у целини описан у
поглављу~\ref{ch:metodologija}.

\section{Хипотезе}
\label{sec:hipoteze}

Полазна теза рада јесте да дизајн мерења представља експерименталну димензију
равноправну избору модела, те да део објављених разлика међу моделима може потицати од
начина припреме података. Да би се та теза могла проверити, рашчлањена је на три
хипотезе:

\begin{itemize}
\item \textbf{Х1:} Варијација у дизајну мерења објашњава удео варијансе губитка
      упоредив са уделом који објашњава избор модела.
\item \textbf{Х2:} Додавање компоненти скокова као улазних променљивих побољшава
      тачност предвиђања. Проверава се на два унапред наведена пара: HAR-J наспрам
      log-HAR и CHAR наспрам log-HAR.
\item \textbf{Х3:} Рангирање модела мења се при преласку са статистичке мере тачности на
      економску евалуацију; мери се Кендаловим коефицијентом $\tau$ између рангирања по
      QLIKE и рангирања по Лопезовој функцији губитка.
\end{itemize}

Централна хипотеза јесте Х1, јер се непосредно односи на постављени циљ рада; Х2
проверава да ли информација о скоковима уопште носи сигнал, без чега налаз о Х1 не би
био интерпретабилан, а Х3 испитује да ли закључак опстаје при промени мерила. Начин
провере сваке хипотезе описан је у одељку~\ref{sec:analiza}. Оба исхода су
информативна: ако се Х1 потврди, следи да дизајн мерења заслужује да буде наведен и
образложен у свакој студији о предвиђању волатилности, једнако као што се наводи
архитектура модела, а ако се не потврди, рад пружа емпиријски основ за супротну
тврдњу, дакле да су реализоване мере довољно робусне да избор међу њима не утиче битно
на поређење модела.

\section{Структура рада}

Рад је организован у пет поглавља. У другом поглављу изложене су теоријске основе:
конструкција реализованих мера и њихово разлагање на континуирану и скоковиту
компоненту, два експериментална фактора која рад укршта, изведене променљиве које чине
улаз модела, као и валидациони протокол и апарат за евалуацију. Треће поглавље описује
методологију: извор и припрему података, дизајн експеримента са фиксираним параметрима,
мере, правила искључивања, план статистичке анализе и ограничења методологије. У
четвртом поглављу приказани су резултати експеримента по нивоима евалуације и
хипотезама, уз дискусију која ограничења из трећег поглавља проширује у светлу
добијених резултата. Пето поглавље сумира главне налазе и наводи правце даљег
истраживања.

\chapter{Теоријске основе}
\label{ch:teorija}

Ово поглавље излаже теоријски оквир истраживања, и то конструкцију реализованих мера, два
експериментална фактора која рад укршта, изведене променљиве које их повезују, те
валидациони протокол и апарат за евалуацију.

\section{Реализована волатилност као изведена величина}

\subsection{Полазни модел цене}

Стандардна поставка у литератури о реализованој волатилности полази од тога да се
промена логаритма цене финансијског инструмента састоји из три дела, и то систематског
кретања (тренда), непрекидних случајних осцилација и повремених скокова
\parencite{abd2007}:
\begin{equation}
dp(t) = \mu(t)\,dt + \sigma(t)\,dW(t) + \kappa(t)\,dq(t)
\label{eq:proces}
\end{equation}
где је $p(t)$ логаритам цене, $\mu(t)$ очекивана стопа промене цене (тренд),
$\sigma(t)$ тренутна (спот) волатилност, $W(t)$ стандардно Брауново кретање, $q(t)$ бројачки процес који
региструје скокове, а $\kappa(t)$ величина скока у тренутку $t$.

Величина од интереса јесте \textbf{квадратна варијација} процеса на дану $t$, која
се разлаже на непрекидни део (интегрисану варијансу) и део који потиче од скокова:
\begin{equation}
\QV_t = \int_{t-1}^{t}\sigma^2(s)\,ds + \sum_{t-1 < s \le t}\kappa^2(s)
\end{equation}

Прва сума представља \textbf{континуирану компоненту}, а друга \textbf{скоковиту
компоненту}; ниједна није непосредно посматрана, већ се обе оцењују из дискретно
узоркованих цена.

\subsection{Реализована варијанса}
\label{subsec:rv}

Ако се дан трговања подели на $n$ интервала једнаке дужине, интрадневни
логаритамски принос на $i$-том интервалу дефинише се као
\begin{equation}
r_{t,i} = p_{t,i} - p_{t,i-1}
\end{equation}

\textbf{Реализована варијанса} (енгл.\ realized variance, RV) јесте сума квадрата тих
приноса:
\begin{equation}
\RV_t = \sum_{i=1}^{n} r_{t,i}^{2}
\end{equation}

Када $n \to \infty$ и у одсуству микроструктурног шума, $\RV_t$ конвергира ка
квадратној варијацији $\QV_t$, дакле ка збиру континуиране и скоковите компоненте.
У раду се доследно користи следећа конвенција: $\RV_t$ је варијанса, а под
волатилношћу се подразумева њен квадратни корен $\sqrt{\RV_t}$; сви модели
предвиђају варијансу, а корен се узима тек при израчунавању вредности при ризику.
У имплементацији се $\RV$ рачуна искључиво из приноса унутар једног дана трговања.
Принос између затварања претходног и отварања текућег дана, у наставку ноћни принос,
никада се не укључује, јер не представља непрекидно трговање и његово укључивање би
контаминирало меру догађајем који припада другој временској скали.

\subsection{Реализована квартичност}

Реализована квартичност (енгл.\ realized quarticity, RQ) оцењује четврти момент и
потребна је као мера прецизности
саме оцене $\RV$, због чега улази у HARQ модел (одељак~\ref{subsec:harq}):
\begin{equation}
\RQ_t = \frac{n}{3}\sum_{i=1}^{n} r_{t,i}^{4}
\end{equation}

Ова величина нестабилна је при малом $n$, што је разлог за апсолутни доњи праг броја
интрадневних опсервација уведен при чишћењу података.

\subsection{Семиваријансе са предзнаком}

Патон и Шепард \parencite{pattonSheppard2015} показују да будућа волатилност зависи
различито од приноса различитог предзнака. Реализоване семиваријансе раздвајају
$\RV$ по предзнаку приноса:
\begin{equation}
\RS_t^{+} = \sum_{i=1}^{n} r_{t,i}^{2}\,\mathbf{1}\{r_{t,i} > 0\},
\qquad
\RS_t^{-} = \sum_{i=1}^{n} r_{t,i}^{2}\,\mathbf{1}\{r_{t,i} < 0\}
\end{equation}

Између њих важи тачан идентитет, који се у имплементацији проверава као тврда
провера исправности:
\begin{equation}
\RS_t^{+} + \RS_t^{-} = \RV_t
\label{eq:semivar-identity}
\end{equation}

\textbf{Варијација скока са предзнаком} дефинише се као њихова разлика:
\begin{equation}
\SJ_t = \RS_t^{+} - \RS_t^{-}
\end{equation}

Дневни принос, потребан за полужни (\textit{leverage}) члан и за евалуацију
вредности при ризику, јесте сума интрадневних приноса, дакле логаритамски принос од
цене прве до цене последње свећице сесије истог дана (од отварања до затварања, без
ноћног приноса):
\begin{equation}
r_{d,t} = \sum_{i=1}^{n} r_{t,i} = p_{t,n} - p_{t,0}
\end{equation}

\section{Први експериментални фактор: дизајн мерења}
\label{sec:dizajn}

Овај одељак описује прву половину факторског експеримента. Три одлуке које се у
литератури најчешће доносе прећутно овде се третирају као експлицитне
експерименталне димензије, означене као А1, А2 и А3.

\subsection{Димензија А1 — фреквенција узорковања}
\label{subsec:a1}

Избор мреже узорковања ставља истраживача пред компромис који нема решење важеће за
све случајеве.

Са статистичке стране, асимптотска варијанса оцене реализоване варијансе опада са
бројем опсервација:
\begin{equation}
\operatorname{Var}(\RV_t) \approx \frac{2}{n}\int_{t-1}^{t}\sigma^{4}(s)\,ds
\end{equation}

Ово сугерише да је чешће узорковање увек боље.

Битно је нагласити да посматрана цена инструмента није ефикасна цена, већ садржи
\textbf{микроструктурни шум}:
\begin{equation}
p_{t,i}^{obs} = p_{t,i}^{*} + u_{t,i}
\end{equation}
где је $p^{*}$ ефикасна (неопсервабилна) цена, а $u$ шум који потиче од распона
између куповне и продајне цене, дискретизације цена и ефеката извршења налога.

Под претпоставком да је шум независан и једнако распоређен, очекивана вредност
реализоване варијансе рачунате из посматраних цена износи
\parencite{zma2005,hansenLunde2006}
\begin{equation}
E\!\left[\RV_t^{obs}\right] = \QV_t + 2n\,E\!\left[u^{2}\right]
\label{eq:noise-bias}
\end{equation}

Други члан расте \textbf{линеарно са $n$}. Пораст фреквенције узорковања стога
смањује варијансу оцене, али истовремено линеарно увећава пристрасност. То је разлог
зашто је петоминутно узорковање постало стандард у литератури. Она представља емпиријски
компромис, а не теоријски оптимум. Лиу, Патон и Шепард \parencite{lps2015} на
широком скупу инструмената налазе да се петоминутна $\RV$ као заменска мера за стварну
варијансу тешко може значајно надмашити.

У овом раду посматрају се три учесталости узорковања цена: једноминутна, петоминутна и
петнаестоминутна. У даљем тексту оне се означавају као нивои $1m$, $5m$ и $15m$ фактора
А1.

Свака мрежа изводи се преузорковањем истих учитаних једноминутних свећица, тако да
димензија А1 мења искључиво агрегацију, а не извор података. Пун дан трговања (390
минута) даје 390, 78 и 26 свећица на мрежама $1m$, $5m$ и $15m$. Дани у којима берза ради
скраћено и трговање се завршава у 13:00, на пример дан после Дана захвалности или
Бадње вече, па трају 210 уместо 390 минута, искључују се из узорка за све три
учесталости узорковања. Тако се сва три нивоа фактора А1 заснивају на истом скупу дана
(одељак~\ref{subsec:ciscenje}).

\subsection{Димензија А2 — естиматори отпорни на скокове}

Скок је нагла, дисконтинуирана промена цене у веома кратком временском интервалу,
најчешће изазвана доласком нове информације на тржиште, на пример објавом резултата
пословања или макроекономских података. У једначини~\eqref{eq:proces} скокове описује
члан $\kappa(t)\,dq(t)$, док непрекидна компонента $\sigma(t)\,dW(t)$ описује
уобичајена, постепена колебања цене.

Циљ сваког естиматора у овој групи јесте да оцени \textbf{искључиво континуирани
део} волатилности, дакле оно што би $\RV$ био да скокова нема. Разлика $\RV$ и
таквог естиматора представља оцену доприноса скокова.

Димензија А2 има три нивоа, односно три естиматора: биповер варијацију (енгл.\ bipower
variation, BPV), медијалну реализовану варијансу (енгл.\ median realized variance,
MedRV) и биповер варијацију са прагом (енгл.\ threshold bipower variation, TrBPV). У
наставку се произвољан естиматор ове групе означава општим симболом $\IVhat_t$ (оцена
интегрисане варијансе, енгл.\ integrated variance, IV). У датом дизајну он узима вредност $\BPV_t$, $\MedRV_t$ или
$\TrBPV_t$ и као такав се користи у формулама тестова и декомпозиције (једначине
\eqref{eq:naive}, \eqref{eq:rj}, \eqref{eq:bns} и \eqref{eq:cond-decomp}).

\paragraph{Биповер варијација (BPV).}
Барндорф-\hspace{0pt}Нилсен и Шепард \parencite{bns2004} уместо квадрата једног
приноса множе
апсолутне вредности два суседна приноса. Пошто скок утиче на само два сабирка и то
без квадрирања, његов утицај је потиснут:
\begin{equation}
\BPV_t = \mu_1^{-2}\cdot\frac{n}{n-1}\sum_{i=2}^{n}\left|r_{t,i}\right|\left|r_{t,i-1}\right|
\end{equation}
где је константа скалирања
\begin{equation}
\mu_1 = E\!\left[|Z|\right] = \sqrt{\frac{2}{\pi}} \approx 0{,}7979,
\qquad
\mu_1^{-2} = \frac{\pi}{2} \approx 1{,}5708
\end{equation}

Фактор корекције за коначан узорак $n/(n-1)$ рачуна се из стварног броја приноса тог
дана након чишћења, никада из фиксиране константе. Дани са рупама у подацима легитимно
носе мањи број свећица, па би фиксиран делилац уносио систематску пристрасност на
сваки такав дан.

\paragraph{Медијална реализована варијанса (MedRV).}
Андерсен, Добрев и Шаумбург \parencite{andersen2012} узимају медијану сваке тројке
узастопних апсолутних приноса, чиме се игноришу и највећи и најмањи члан тројке:
\begin{equation}
\MedRV_t = \frac{\pi}{6 - 4\sqrt{3} + \pi}\cdot\frac{n}{n-2}
\sum_{i=2}^{n-1}\med\!\left(|r_{t,i-1}|,\,|r_{t,i}|,\,|r_{t,i+1}|\right)^{2}
\end{equation}

Овај естиматор отпорнији је од BPV-а у ситуацији када се скокови јаве у узастопним
приносима, дакле у случају у коме је BPV-ов производ двају суседних приноса потпуно
контаминиран.

\paragraph{Биповер варијација са прагом (TrBPV).}
Трећа варијанта, по угледу на Корсија, Пирина и Рена \parencite{corsiPirinoReno2010},
уклања приносе изнад прага изведеног из робусне оцене волатилности пре него што се
формирају унакрсни производи:
\begin{equation}
\TrBPV_t = \mu_1^{-2}\cdot\frac{n}{n-1}
\sum_{i=2}^{n}\left|\tilde{r}_{t,i}\right|\left|\tilde{r}_{t,i-1}\right|
\end{equation}
где је принос са прагом дефинисан као
\begin{equation}
\tilde{r}_{t,i} = r_{t,i}\cdot\mathbf{1}\!\left\{|r_{t,i}| \le \vartheta\hat{\sigma}_t\right\}
\end{equation}

Дневна робусна скала оцењује се једном за цео дан, преко медијане апсолутног приноса
скалиране на јединицу стандардне девијације под нормалношћу:
\begin{equation}
\hat{\sigma}_t = \frac{\med\left(|r_{t,i}|\right)}{0{,}6745}
\end{equation}

Праг је фиксиран на $\vartheta = 3{,}0$ кроз цео простор дизајна. Ова одлука је
намерна, будући да би увођење прага као четврте димензије дизајна увећало мрежу са 36 на 108
комбинација, а пошто се параметар односи само на један од три естиматора, дизајн би
постао неуравнотежен и удели варијансе по фактору у анализи варијансе (енгл.\ analysis
of variance, ANOVA) постали би
неинтерпретабилни.

\subsection{Димензија А3 — тестови за детекцију скокова}
\label{subsec:a3}

\paragraph{Наивни праг.}
Дан се проглашава даном са скоком ако апсолутна вредност укупног дневног приноса
пређе фиксиран број стандардних девијација, дакле тест је двостран:
\begin{equation}
S_t^{naive} = \frac{\left|r_{d,t}\right|}{\sqrt{\IVhat_t}},
\qquad
\text{скок ако } S_t^{naive} > 3
\label{eq:naive}
\end{equation}

Кључан детаљ имплементације јесте да се стандардна девијација у имениоцу изводи из
естиматора отпорног на скокове, никада из стандардне девијације рачунате преко целог
дана. Друга варијанта укључила би сам скок који се тестира у сопствени именилац, па
би довољно велик скок надувао сопствени праг и избегао детекцију. Овај тест не
користи се као озбиљна метода, већ као доња граница квалитета.

\paragraph{BNS тест.}
Тест Барндорф-Нилсена и Шепарда (енгл.\ Barndorff-Nielsen--Shephard test, BNS)
\parencite{bns2006} пореди $\RV$ и естиматор отпоран на скокове, скалирано триповер
квартичношћу (енгл.\ tripower quarticity, TPQ); у облику односа тест су проучавали
Хуанг и Таухен \parencite{huangTauchen2005}. Релативни допринос скокова дефинише се као
\begin{equation}
RJ_t = \frac{\RV_t - \IVhat_t}{\RV_t}
\label{eq:rj}
\end{equation}

Триповер квартичност, потребна за нормализацију, износи
\begin{equation}
\TPQ_t = n\,\mu_{4/3}^{-3}\cdot\frac{n}{n-2}
\sum_{i=3}^{n}\left|r_{t,i}\right|^{4/3}\left|r_{t,i-1}\right|^{4/3}\left|r_{t,i-2}\right|^{4/3}
\end{equation}

Тест статистика има асимптотску стандардну нормалну расподелу под нултом хипотезом о
одсуству скокова:
\begin{equation}
Z_t = \frac{RJ_t}{\sqrt{\theta\cdot\dfrac{\TPQ_t}{n\,\IVhat_t^{2}}}}\ \sim\ N(0,1)
\label{eq:bns}
\end{equation}

Тест је једностран, јер скок може само да увећа $\RV$ у односу на $\IVhat$. Дан се
проглашава даном са скоком ако је $p = 1 - \Phi(Z_t) < \alpha$, уз $\alpha = 0{,}001$.
Константа асимптотске варијансе $\theta$ зависи од естиматора у имениоцу. За $\BPV$
и $\TrBPV$ она износи
\begin{equation}
\theta = \left(\frac{\pi}{2}\right)^{2} + \pi - 5 \approx 0{,}609,
\end{equation}
а за $\MedRV$ износи $\theta \approx 0{,}96$ \parencite{andersen2012}.

Нормализација са $n\,\IVhat_t^{2}$, а не са $\RV_t^{2}$, јесте оно што даје исправну
асимптотску скалу.

% Дуг наслов пасуса и цитат у загради одмах иза њега не могу да стану у први
% ред, а hyperref их пакује у кутију па се не могу преломити; \emergencystretch
% дозвољава TeX-у да тај ред мало растегне уместо да га пусти преко маргине.
\begingroup
\emergencystretch=2em
\paragraph{Тест Лија и Микланда (енгл.\ Lee--Mykland test, LM).}
Ли и Микланд \parencite{leeMykland2008} не тестирају дан као целину, већ сваки
\textbf{интрадневни принос} посебно, делећи га локалном оценом волатилности
изведеном из клизног прозора од $K$ претходних опсервација:
\begin{equation}
L_{t,i} = \frac{r_{t,i}}{\hat{\sigma}_{t,i}}
\end{equation}
\par
\endgroup

Локална оцена волатилности рачуна се биповер поступком без корекције фактором
$\mu_1^{-2}$, као у изворном раду Лија и Микланда:
\begin{equation}
\hat{\sigma}_{t,i}^{2} = \frac{1}{K-1}\sum_{j=i-K+1}^{i-1}\left|r_j\right|\left|r_{j-1}\right|
\label{eq:local-sigma}
\end{equation}

Изостављена корекција значи да статистика под нултом хипотезом није стандардно
нормална, па се критична вредност изводи уз ту конвенцију. Максимум апсолутних
статистика дана приближно прати Гумбелову расподелу, са параметрима положаја и размере
\begin{equation}
C_n = \frac{\sqrt{2\ln n}}{c} - \frac{\ln\pi + \ln(\ln n)}{2c\sqrt{2\ln n}},
\qquad
S_n = \frac{1}{c\sqrt{2\ln n}}
\end{equation}
где је $c = \mu_1 = \sqrt{2/\pi}$. Критична вредност на нивоу $0{,}999$ износи
\begin{equation}
q_n = C_n + S_n\beta^{*},
\qquad
\beta^{*} = -\ln\!\left(-\ln 0{,}999\right)
\end{equation}

Дан се проглашава даном са скоком ако важи $\max_i |L_{t,i}| > q_n$. Критична
вредност зависи од $n$ и рачуна се за сваки дан посебно. Фиксиран праг (на пример
$|L| > 4{,}0$) помешао би димензију А1 са димензијом А3, јер би исти праг значио различит
ниво значајности на различитим мрежама.

Дужина локалног прозора $K$ фиксирана је на \textbf{четири сесије трговања на свакој
мрежи} ($K = 1560/312/104$ за мреже $1m/5m/15m$). Ово је намерна одлука, а не
преузимање вредности из литературе. Када би се дужина прозора изражена у сесијама
мењала између мрежа, класификација скока зависила би од две ствари које се обе мењају
са димензијом А1, што би искривило удео варијансе приписан дизајну мерења.

\paragraph{LM тест са контролом стопе лажних открића (LM-FDR).}
Гумбелова критична вредност представља контролу грешке на нивоу породице
\textbf{унутар једног дана}, преко свих $n$ интрадневних тестова тог дана, али ништа
не говори о томе шта се дешава \textbf{преко дана}. Четврта варијанта теста зато мења
критеријум, па се уместо фиксиране критичне вредности по дану контролише стопа лажних
открића (енгл.\ false discovery rate, FDR), односно очекивани удео лажних међу
проглашеним данима са скоком, преко клизног прозора дана.

Сваком дану најпре се додељује дневна $p$-\hspace{0pt}вредност из Гумбелове
расподеле максимума:
\begin{equation}
p_t = 1 - \exp\!\left[-\exp\!\left(-\frac{\max_i|L_{t,i}| - C_n}{S_n}\right)\right]
\end{equation}

Бенјамини–\hspace{0pt}Хохбергова процедура (енгл.\ Benjamini--\hspace{0pt}Hochberg
procedure, BH) \parencite{benjaminiHochberg1995} примењује се затим
преко клизног прозора од $W$ претходних дана плус текућег дана. За сортиране
$p$-\hspace{0pt}вредности $p_{(1)} \le p_{(2)} \le \dots \le p_{(m)}$ проналази се највећи индекс
који задовољава
\begin{equation}
k^{*} = \max\left\{k\ :\ p_{(k)} \le \frac{k}{m}q\right\},
\qquad q = 0{,}05
\end{equation}
а текући дан проглашава се даном са скоком ако је $p_t \le p_{(k^{*})}$. Прозор износи
$W = 250$ дана, што одговара приближно једној години трговања. Клизни прозор је
обавезан, јер би примена преко целог узорка класификовала дан користећи $p$-вредности из
година које му тек следе, што представља гледање унапред.

Два критеријума контролишу различите величине: Гумбелов праг вероватноћу бар једне
лажне детекције унутар дана, а BH процедура удео лажних међу
проглашеним данима. Пошто BH праг зависи од расподеле $p$-вредности у прозору, при
овим нивоима он је блажи од Гумбеловог, па LM-FDR означава више дана него обичан LM
тест. Измерени удели дати су у одељку~\ref{sec:rez-uzorak}.

\paragraph{Период загревања.}
Клизни прозори нису дефинисани на почетку серије, будући да локална оцена волатилности захтева
$K$ сесија, прозор за контролу стопе лажних открића 250 дана (до тада се користи обичан
LM тест), а месечни агрегат 22 дана. Сви ти периоди падају унутар 2010–2013. године,
која се никада не вреднује, па загревање не утиче на извештаване резултате.

\subsection{Условна декомпозиција на континуирану и скоковиту компоненту}
\label{subsec:decompose}

Ово је место на коме се димензије А2 и А3 срећу и једино место у ком се обе одлуке
истовремено уграђују у улазне податке модела.

\textbf{Безусловна} декомпозиција, $J_t = \max(\RV_t - \IVhat_t, 0)$, зависи само од
естиматора, па би заменом теста вредности $J$ и $C$ остале идентичне, а димензија А3
била безефектна за сваки улаз модела.

\textbf{Условна} декомпозиција \parencite{abd2007} то исправља тако што скоковиту
компоненту признаје само на данима на којима је тест значајан:
\begin{equation}
J_t = \mathbf{1}\{\text{тест значајан на дан } t\}\cdot\max\!\left(\RV_t - \IVhat_t,\ 0\right)
\label{eq:cond-decomp}
\end{equation}
\begin{equation}
C_t = \RV_t - J_t
\end{equation}

Сада $J$ зависи и од тога који је естиматор употребљен (димензија А2) и од тога који је
тест прогласио дан даном са скоком (димензија А3). На данима без значајног теста важи
$J_t = 0$ и сва варијација приписује се континуираном делу, што је и теоријски
исправно, будући да је просечан јаз $(\RV - \IVhat)/\RV$ на данима на којима тест није
значајан мали, што значи да је сирови јаз на тим данима шум естиматора, а не скок. На
референтном дизајну коначног панела тај јаз износи $0{,}060$ на данима без значајног
теста, наспрам $0{,}388$ на данима са значајним тестом. Сирови јаз на неозначеним данима
јесте дакле реда величине шума, док на означеним данима носи величину која се не може
приписати шуму естиматора.

Иста формула користи се за све четири варијанте теста, јер би другачија формула за
једну варијанту помешала ефекат теста са ефектом формуле у истом фактору анализе
варијансе.

\subsection{Простор дизајна мерења}

Комбиновањем трију фактора добија се простор дизајна
\begin{equation}
D = A1 \times A2 \times A3,
\qquad
|D| = 3 \times 3 \times 4 = 36
\end{equation}

Сваки дизајн идентификује се кратким низом који спаја ознаке мреже, естиматора и
теста, на пример \texttt{5m\_\_bpv\_\_bns}. Тај низ касније постаје кључ фактора у
декомпозицији варијансе.

\subsection{Изведене променљиве}
\label{sec:izvedene}

Изведене променљиве настају из дизајна мерења и заједно чине улаз предиктивних модела.
Пошто свака од њих зависи од једне или више димензија описаних у претходним
пододељцима, промена дизајна мења њихове вредности, а тиме и све што модели виде.

Из дневних мера граде се недељни и месечни просеци, који чине улаз HAR породице
модела:
\begin{equation}
\RV_{w,t} = \frac{1}{5}\sum_{j=0}^{4}\RV_{t-j},
\qquad
\RV_{m,t} = \frac{1}{22}\sum_{j=0}^{21}\RV_{t-j}
\end{equation}

Исти поступак примењује се и на компоненте $C$ и $J$, дајући величине
$C_d, C_w, C_m$ и $J_d, J_w, J_m$. Три временска хоризонта одговарају
хипотези о хетерогеном тржишту \parencite{muller1997,corsi2009}, према којој
учесници реагују различитим брзинама.

Изведене променљиве приказане су у табели~\ref{tab:izvedene}, заједно са
назнаком од којих димензија дизајна мерења свака зависи. Та зависност је
суштинска за тумачење резултата. Променљиве које зависе од димензија А2 и А3
мењају вредност са сваком изменом дизајна, док оне које зависе само од А1
остају идентичне кроз све комбинације естиматора и теста.

\begin{table}[htbp]
\centering
\caption{Изведене променљиве и њихова зависност од димензија дизајна мерења}
\label{tab:izvedene}
\begin{tabularx}{\textwidth}{L{3.4cm} X L{2.2cm}}
\toprule
Група & Променљиве & Зависи од \\
\midrule
HAR компоненте &
$\RV_d$, $\RV_w$, $\RV_m$ &
А1 \\
\addlinespace
Компоненте скокова &
$C_d$, $C_w$, $C_m$, $J_d$, $J_w$, $J_m$, бинарни индикатор скока &
А1, А2, А3 \\
\addlinespace
Семиваријансе &
$\RS^{+}$, $\RS^{-}$, варијација са предзнаком $\SJ$ &
А1 \\
\addlinespace
Мера прецизности &
реализована квартичност $\RQ$ &
А1 \\
\addlinespace
Полужни члан &
дневни принос $r_d$ &
А1 \\
\addlinespace
Егзогена променљива &
$\VIX$ са кашњењем од једног дана &
— \\
\bottomrule
\end{tabularx}
\end{table}

Основ за разлагање на семиваријансе јесте налаз Патона и Шепарда
\parencite{pattonSheppard2015} да је негативна семиваријанса знатно значајнија
за будућу волатилност од позитивне. Индекс имплицитне волатилности узима се
искључиво са кашњењем од једног дана, јер би коришћење вредности из истог дана за
предвиђање тог истог дана представљало гледање унапред.

\section[Други експериментални фактор: архитектуре предиктивних модела]{Други експериментални фактор:\\ архитектуре предиктивних модела}
\label{sec:modeli}

\subsection{Конструкција циљне променљиве}
\label{subsec:target}

Циљна величина јесте \textbf{$h$-дневни просек} будуће реализоване варијансе, што је
Корсијева изворна HAR конструкција и оно што свака студија у овој области извештава:
\begin{equation}
y_{t,h} = \frac{1}{h}\sum_{i=1}^{h}\RV_{t+i}
\label{eq:target}
\end{equation}

Хоризонти су $h \in \{1, 5, 22\}$, што одговара дану, недељи и месецу. Циљ се моделује
у логаритамском простору:
\begin{equation}
\ln y_{t,h} = \ln\!\left(\frac{1}{h}\sum_{i=1}^{h}\RV_{t+i}\right)
\end{equation}

Разлог за логаритамску трансформацију јесте изражена десна асиметрија расподеле
реализоване волатилности. Мали број модела, чија је позитивност структурна, уместо
тога користи циљ у нивоима (одељци \ref{subsec:pozitivnost} и \ref{subsec:mem}).

\subsection{Проблем повратне трансформације и Јенсенова неједнакост}
\label{subsec:jensen}

Ако модел предвиђа $\ln y$, наивна повратна трансформација
$\exp(\widehat{\ln y})$ \textbf{систематски потцењује} ниво. Разлог је Јенсенова
неједнакост, по којој за конвексну функцију важи
\begin{equation}
E\!\left[\exp(X)\right] \ge \exp\!\left(E[X]\right)
\end{equation}

Величина $\exp(\widehat{\ln y})$ представља оцену условне \textbf{медијане} само ако
је расподела грешке у логаритамском простору симетрична; у општем случају она је
оцена нивоа који одговара средини логаритама, а не условне средине. Пошто су QLIKE и
корен средње квадратне грешке (енгл.\ root mean squared error, RMSE) минимизовани
условном средином, потребна је
корекција.

Уобичајена Гаусова корекција гласила би $a = \tfrac{1}{2}\sigma^{2}$, али она
претпоставља да су резидуали тачно нормално распоређени. Резидуали реализоване
волатилности то не морају бити, посебно код модела заснованих на стаблима. У овом
раду корекција је стога \textbf{непараметарска} и одговара Дуановом поступку
размазивања \parencite{duan1983}:
\begin{equation}
a = \ln\!\left(\frac{1}{m}\sum_{k=1}^{m}\exp(e_k)\right),
\qquad
e_k = \ln y_k - \widehat{\ln y}_k
\label{eq:correction}
\end{equation}

Повратна трансформација гласи
\begin{equation}
\hat{y} = \exp\!\left(\widehat{\ln y} + a\right)
\end{equation}

Корекција $a$ оцењује се на клизном прозору ранијих ван-узорачких грешака, како је
описано у одељку~\ref{subsec:korekcija-prozor}. Никада се не оцењује на резидуалима
унутар узорка, који су код модела заснованих на стаблима близу нуле.

За моделе који предвиђају \textbf{збир компоненти}, корекција мора бити оцењена на
већ сабраној прогнози, и то као мултипликативни фактор облика однос сума:
\begin{equation}
m^{*} = \frac{\sum_{k=1}^{N} y_k}{\sum_{k=1}^{N}\hat{y}_k^{comb}},
\qquad
\hat{y}^{corr} = m^{*}\hat{y}^{comb}
\end{equation}

Примена засебне корекције по компоненти и сабирање након корекције оставило би збир
без икакве контролисане корекције, јер Јенсенова неједнакост није адитивна преко
компоненти.

\subsection{Механизми позитивности}
\label{subsec:pozitivnost}

Волатилност је по дефиницији ненегативна, па сваки модел мора гарантовати ненегативну
прогнозу. У овом раду постоје \textbf{четири намерно различита механизма}, који се
никада не обједињују, управо зато што рад испитује да ли сам механизам позитивности, а
не само скуп улазних променљивих, утиче на тачност. У табели~\ref{tab:pozitivnost}
приказани су сви механизми са припадајућим моделима; први механизам обухвата и
варијанту у којој се једна корекција примењује на већ сабрану прогнозу.

\begin{table}[htbp]
\centering
\caption{Механизми обезбеђивања ненегативне прогнозе}
\label{tab:pozitivnost}
\begin{tabularx}{\textwidth}{L{3.4cm} L{4.6cm} X}
\toprule
Механизам & Модели & Начин гарантовања \\
\midrule
Логаритам и непараметарска корекција &
\texttt{naive}, \texttt{log\_har}, \texttt{char}, \texttt{har\_j}, \texttt{shar},
\texttt{harq}, \texttt{har\_iv}, \texttt{arfima}, \texttt{lightgbm},
\texttt{xgboost}, \texttt{ttm\_rv}, \texttt{ttm\_c} &
Оцењивање у логаритамском простору, повратна трансформација преко експоненцијалне
функције уз корекцију из једначине~\eqref{eq:correction} \\
\addlinespace
Логаритам и једна корекција на сабраној прогнози &
\texttt{ttm\_c\_plus\_j}, \texttt{log\_har\_ttm} &
Компоненте се сабирају у нивоима (уз одсецање на нули за $J$), а на збир се примењује
један мултипликативни фактор облика однос сума \\
\addlinespace
Ненегативни најмањи квадрати у нивоима &
\texttt{har} &
Ненегативни коефицијенти уз ненегативне регресоре \\
\addlinespace
Структурна мултипликативна форма &
\texttt{mem} &
Линеарна комбинација са ненегативним коефицијентима; $\mu_t > 0$ јер је $\omega > 0$ \\
\bottomrule
\end{tabularx}
\end{table}

\subsection{Наивна основа}

Наивни модел преноси последњу посматрану вредност као равну прогнозу за цео хоризонт,
и то у логаритамском простору:
\begin{equation}
\widehat{\ln y}_{t,h} = \ln\RV_{d,t}
\end{equation}

Он не садржи ниједан оцењен параметар, али пролази кроз исту повратну трансформацију
и исту непараметарску корекцију као сваки други логаритамски модел
(одељак~\ref{subsec:jensen}), јер би поређење некориговане са коригованим прогнозама
мерило корекцију, а не модел. Служи као доња граница коју сваки озбиљан модел мора
надмашити.

\subsection{HAR модел и његове варијанте}
\label{subsec:har}

\paragraph{Полазна интуиција.}
Корси \parencite{corsi2009} полази од хипотезе о хетерогеном тржишту
\parencite{muller1997}, по којој учесници реагују различитим брзинама (дневни трговци, недељни
портфолио менаџери и дугорочни институционални инвеститори), па волатилност
представља каскаду компоненти различитих временских хоризоната. Дуготрајно памћење
тако се апроксимира без формалног дугомеморијског процеса.

\needspace{5\baselineskip}
\paragraph{Логаритамски HAR.}
Стандардна форма и референтни модел целог рада:
\nopagebreak
\begin{equation}
\ln y_{t,h} = \beta_0 + \beta_d \ln\RV_{d,t} + \beta_w \ln\RV_{w,t} + \beta_m \ln\RV_{m,t} + \varepsilon_t
\label{eq:loghar}
\end{equation}

Оцењује се обичним методом најмањих квадрата, а прогноза се враћа у ниво преко
непараметарске корекције из одељка~\ref{subsec:jensen}. Овај модел, а не
\texttt{har}, јесте основа за релативни QLIKE.

\paragraph{HAR у нивоима са ненегативним најмањим квадратима.}
Иста скупина улазних променљивих, али потпуно другачији механизам позитивности:
\begin{equation}
y_{t,h} = \beta_0 + \beta_d \RV_{d,t} + \beta_w \RV_{w,t} + \beta_m \RV_{m,t} + \varepsilon_t
\end{equation}
уз ограничење да су сви коефицијенти ненегативни:
\begin{equation}
\hat{\beta} = \arg\min_{\beta \ge 0}\left\|y - X\beta\right\|^{2}
\end{equation}

Пошто су сви регресори ненегативни (реч је о варијансама) и сви коефицијенти
ненегативни, прогноза не може бити негативна по конструкцији, без логаритамске
трансформације и без корекције.

\paragraph{CHAR.}
Уместо $\RV$, као регресор користи се искључиво континуирана компонента:
\begin{equation}
\ln y_{t,h} = \beta_0 + \beta_d \ln C_{d,t} + \beta_w \ln C_{w,t} + \beta_m \ln C_{m,t} + \varepsilon_t
\end{equation}

\paragraph{HAR-J.}
Континуирана и скоковита компонента улазе као \textbf{одвојени} регресори, чиме се
тестира да ли скокови имају другачију перзистентност \parencite{abd2007}. Пошто је
$J_t = 0$ на већини дана, скоковита компонента улази кроз $\ln(1 + J_{d,t})$, што је
дефинисано и при $J = 0$:
\begin{equation}
\ln y_{t,h} = \beta_0 + \beta_d \ln C_{d,t} + \beta_w \ln C_{w,t} + \beta_m \ln C_{m,t}
+ \beta_J \ln\!\left(1 + J_{d,t}\right) + \varepsilon_t
\end{equation}

\paragraph{SHAR.}
Дневна компонента раздваја се по предзнаку, чиме се хвата асиметрија коју описују
Патон и Шепард \parencite{pattonSheppard2015}:
\begin{equation}
\ln y_{t,h} = \beta_0 + \beta^{+}\ln\RS_{t}^{+} + \beta^{-}\ln\RS_{t}^{-}
+ \beta_w \ln\RV_{w,t} + \beta_m \ln\RV_{m,t} + \varepsilon_t
\end{equation}

\paragraph{HAR-IV.}
HAR проширен индексом имплицитне волатилности са кашњењем од једног дана, чиме се
тестира да ли тржиште опција садржи информацију које нема у историји цена:
\begin{equation}
\ln y_{t,h} = \beta_0 + \beta_d \ln\RV_{d,t} + \beta_w \ln\RV_{w,t} + \beta_m \ln\RV_{m,t}
+ \beta_{IV}\VIX_{t-1} + \varepsilon_t
\end{equation}

Индекс улази у нивоима, као у изворној HAR-IV спецификацији; реч је о унапред
фиксираној одлуци о моделу, која се не мења кроз простор дизајна.

\subsection{HARQ — модел изграђен на грешци мерења}
\label{subsec:harq}

HARQ \parencite{bpq2016} теоријски је најближи тези овог рада. Идеја је да коефицијент
уз дневну компоненту не буде константан, већ да зависи од тога колико је мерење тог
дана било поуздано. У овом раду користи се варијанта у логаритамском простору, док је
изворни модел постављен у нивоима:
\begin{equation}
\ln y_{t,h} = \beta_0 + \left(\beta_d + \beta_Q z_t\right)\ln\RV_{d,t}
+ \beta_w \ln\RV_{w,t} + \beta_m \ln\RV_{m,t} + \varepsilon_t
\label{eq:harq}
\end{equation}
где је $z_t$ стандардизована вредност квадратног корена реализоване квартичности:
\begin{equation}
z_t = \frac{\sqrt{\RQ_t} - \overline{\sqrt{\RQ}}}{s_{\sqrt{\RQ}}}
\end{equation}

Средина и стандардна девијација рачунају се на подацима за обуку текућег оцењивања и
замрзавају се до наредног, никада се не оцењују на редовима који се предвиђају.
Стандардизација је нужна: $\sqrt{\RQ}$ је реда $10^{-4}$, док је $\ln\RV_d$ реда
$-9$, па би сиров производ био занемарљив у односу на основни регресор. Ако је оцењено $\beta_Q < 0$, модел смањује тежину дневне компоненте на
данима на којима је мерење било зашумљено; предзнак је емпиријско питање, а не
претпоставка модела.

Значај овог модела за рад је концептуалан. Он представља постојећи, у литератури
признат модел изграђен управо на претпоставци да грешка мерења систематски утиче на
прогнозу. Овај рад ту идеју проширује тако што грешку мерења не моделује унутар једног
модела, већ је варира експериментално.

\subsection[ARFIMA — експлицитно дугорочно памћење]{ARFIMA (енгл.\ autoregressive
fractionally integrated moving average) — експлицитно дугорочно памћење}

За разлику од HAR-а, који дуготрајно памћење апроксимира каскадом, ARFIMA га моделује
експлицитно преко фракционог диференцирања:
\begin{equation}
(1-L)^{d}\left(\ln\RV_t - \mu\right) = \Psi(L)\varepsilon_t
\end{equation}

Параметар фракционог диференцирања $d$ оцењује се локалним Витловим естиматором
\parencite{robinson1995}, полупараметарским поступком који користи само ниске
фреквенције периодограма, уз ширину појаса $m = \lfloor n^{0{,}65}\rfloor$; редови
$(p,q)$ бирају се по Бајесовом информационом критеријуму преко скупа $\{0,1,2\}^{2}$, а
прогноза се ре-интегрише у нивое инверзним оператором $(1-L)^{-d}$.

Модел се оцењује \textbf{независно по симболу}, за разлику од HAR породице која се
може обједињавати кроз попречни пресек.

\subsection[Мултипликативни модел грешке (MEM)]{Мултипликативни модел грешке
(енгл.\ multiplicative error model, MEM)}
\label{subsec:mem}

Енгл \parencite{engle2002} предлаже мултипликативну поставку у којој се реализована
варијанса представља као производ условне средине и позитивног случајног члана:
\begin{equation}
\RV_t = \mu_t\varepsilon_t,
\qquad \varepsilon_t > 0,
\qquad E[\varepsilon_t] = 1
\end{equation}

Условна средина прати рекурзију у стилу уопштеног ауторегресивног условно
хетероскедастичног модела (енгл.\ generalized autoregressive conditional
heteroskedasticity, GARCH):
\begin{equation}
\mu_t = \omega + \alpha\RV_{t-1} + \beta\mu_{t-1}
\label{eq:mem}
\end{equation}

\textbf{Позитивност је структурна}: $\mu_t$ је линеарна комбинација ненегативних
величина са ненегативним коефицијентима, а строго је позитивна јер је $\omega > 0$,
без икакве логаритамске трансформације и без корекције повратне трансформације.

Оцењивање се врши квази-методом максималне веродостојности уз експоненцијалну
веродостојност:
\begin{equation}
\ell(\omega,\alpha,\beta) = \sum_{t}\left[-\ln\mu_t - \frac{\RV_t}{\mu_t}\right]
\end{equation}
уз ограничења $\omega > 0$, $\alpha \ge 0$, $\beta \ge 0$ и услов стационарности
$\alpha + \beta < 1$. Разлог за избор експоненцијалне веродостојности јесте робусност:
оцена остаје конзистентна за било коју стварну расподелу грешке са позитивним носачем
и јединичном средином, што је мултипликативни аналог резултата Болерслева и Вулдриџа
\parencite{bollerslevWooldridge1992} за GARCH моделе.

Прогноза за хоризонт $h$ добија се пројекцијом рекурзије унапред, при чему за $k > 1$
важи $E[\mu_{t+k}\mid F_t] = \omega + (\alpha+\beta)E[\mu_{t+k-1}\mid F_t]$, а коначна
прогноза јесте просек тих вредности преко $k = 1,\dots,h$, што одговара конструкцији
циља из одељка~\ref{subsec:target}. И овај модел оцењује се независно по
симболу.

\subsection{Модели засновани на стаблима}
\label{subsec:stabla}

\paragraph{Градијентно појачавање.}
LightGBM \parencite{ke2017} и XGBoost \parencite{chenGuestrin2016} граде адитивни
ансамбл плитких стабала одлучивања, у коме свако наредно стабло учи на негативном
градијенту функције губитка претходних, уз регуларизацију по броју листова и
вредностима у њима.

Предност ових модела јесте способност хватања нелинеарних интеракција које линеарни
HAR не може представити. Улазни скуп представља унију свих променљивих HAR породице,
проширену полужним чланом, при чему улазе у истом облику у коме их виде и линеарни
модели, и то реализоване мере у логаритму, скоковита компонента кроз $\ln(1 + J_d)$, а
дневни принос и индекс имплицитне волатилности у нивоу:
\begin{equation}
\begin{split}
X = \big\{& \ln\RV_d, \ln\RV_w, \ln\RV_m, \ln C_d, \ln C_w, \ln C_m,
\ln(1 + J_d), \\
& \ln\RS^{+}, \ln\RS^{-}, \sqrt{\RQ}\cdot\ln\RV_d, \VIX_{t-1}, r_d\big\}
\end{split}
\end{equation}

Интеракциони члан HARQ модела улази нестандардизован, јер је стандардизација питање
линеарне спецификације, а стабла поред њега виде и сам $\ln\RV_d$.

Хиперпараметри су \textbf{фиксирани и наведени}, а не подешавани по прозору (200
стабала, максимална дубина 4, 15 листова, стопа учења $0{,}05$). Разлог је
методолошки, јер би подешавани хиперпараметри учинили да допринос варијанси одражава
уложен труд у претрагу, а не саму архитектуру, што би обесмислило поређење са осталим
факторима.

\subsection{Претходно обучени модели опште намене}
\label{subsec:ttm}

Најновији правац чине модели опште намене за временске серије (енгл.
\textit{foundation models}), дакле модели претходно обучени на великом броју разноврсних
временских серија, који се потом примењују \textbf{без икаквог додатног обучавања}
(\textit{zero-shot}).

У овом раду користи се IBM Granite TinyTimeMixer (у наставку TTM)
\parencite{ekambaram2024}, са
\textbf{фиксираном верзијом} контролне тачке, дакле модел
\url{ibm-granite/granite-timeseries-ttm-r2}, издање \texttt{512-192-r2}, са дужином
контекста 512 и дужином излаза 192 корака. Фиксирање издања је методолошки неопходно, јер су различита издања обучена
на различитим скуповима података, па би провера контаминације над нефиксираним
моделом била бесмислена. Упркос имену, реч је о архитектури типа MLP-Mixer, а не о
трансформеру, па се временски и канални правац обрађују наизменичним потпуно повезаним
слојевима, без механизма пажње.

Модел се примењује \textbf{униваријантно, по симболу}, па је улаз последњих 512
вредности серије тог симбола, а ако серија на почетку узорка има мање од 512
опсервација, недостајући део попуњава се понављањем прве расположиве вредности (лево
допуњавање). Модел враћа 192 корака унапред, а прогноза за хоризонт $h$ јесте просек
првих $h$ корака, што одговара конструкцији циља из одељка~\ref{subsec:target}.

Реализоване варијанте приказане су у табели~\ref{tab:ttm}.

\begin{table}[htbp]
\centering
\caption{Варијанте претходно обученог модела и њихова зависност од фактора мерења}
\label{tab:ttm}
\begin{tabular}{lll}
\toprule
Варијанта & Улаз & Зависи од \\
\midrule
\texttt{ttm\_rv}        & серија $\ln \RV_d$        & мреже \\
\texttt{ttm\_c}         & серија $\ln C_d$          & мреже, естиматора, теста \\
\texttt{ttm\_c\_plus\_j} & серије $C$ и $J$ одвојено & мреже, естиматора, теста \\
\texttt{log\_har\_ttm}  & хибрид \texttt{log\_har} и \texttt{ttm\_rv} & мреже \\
\bottomrule
\end{tabular}
\end{table}

Модел \texttt{ttm\_c\_plus\_j} прогнозира континуирану компоненту на уобичајеном
логаритамском путу, а скоковиту компоненту \textbf{у нивоима}, јер $J_t = 0$ важи на
већини дана, па је $\ln J$ недефинисан. Обе прогнозе одсецају се на нули, сабирају у
нивоима, а на збир се примењује \textbf{једна} комбинована корекција из
одељка~\ref{subsec:jensen}.

\needspace{6\baselineskip}
Хибридни модел комбинује прогнозе модела log-HAR и TTM-RV у логаритамском простору,
пре корекције, са тежинама \textbf{фиксираним на $0{,}5$ и $0{,}5$ и никада
оптимизованим}:
\nopagebreak
\begin{equation}
\widehat{\ln y}^{\,hybrid} = 0{,}5\,\widehat{\ln y}^{\,log\_har} + 0{,}5\,\widehat{\ln y}^{\,ttm\_rv}
\end{equation}

Тежине нису оптимизоване из истог разлога из ког ни хиперпараметри стабала нису
подешавани (одељак~\ref{subsec:stabla}).

Пошто датум до кога сежу подаци за обуку пада унутар периода евалуације,
пре сваке прогнозе овог модела спроводи се провера преклапања, описана у
одељку~\ref{subsec:kontaminacija}.

\section{Валидациони протокол}

\subsection{Проширујући прозор}
\label{subsec:validacija}

Модели се оцењују на \textbf{проширујућем прозору} са фиксним почетком. Почетак скупа за
обуку је 1. јануар 2010. године и никада се не помера, док се крај скупа за обуку
помера унапред како време одмиче. Прозор за обуку никада није клизни, јер фиксан
почетак обезбеђује да се сав расположиви историјски узорак увек користи.

Уместо поделе на неколико великих прозора, модели се поново оцењују
\textbf{сваких 21 дан трговања}, што одговара приближно једном месецу. Свако
оцењивање користи искључиво податке који претходе дану прогнозе:
\begin{equation}
\text{обука за дан } t:\quad
\left[\,T_0,\ \tau(t)\,\right),
\qquad
T_0 = \text{2010-01-01},
\qquad
\tau(t) \le t
\end{equation}
где је $\tau(t)$ дан последњег оцењивања пре дана $t$, дакле почетак интервала је
укључен, а крај искључен. Из скупа за обуку притом се изостављају и они редови чија
циљна вредност још није остварена до дана оцењивања. На хоризонту $h$ последњих $h$
дана пре тог датума има циљ који се тек формира, па би њихово укључивање значило
оцењивање на подацима који у том тренутку нису били познати. На хоризонту $h = 22$ то
је нетривијалан део краја скупа за обуку.

Између два оцењивања параметри модела остају замрзнути. Међутим, улазне величине и
стање модела ажурирају се свакодневно, при чему се HAR агрегати, компоненте,
семиваријансе и рекурзивно стање мултипликативног модела рачунају из података
доступних закључно са даном $t$. Тиме модел у сваком тренутку користи
најсвежије податке као улаз, али никада параметре оцењене на будућности.

Претходно обучени модели опште намене не оцењују се уопште, него се примењују без
додатног обучавања, па се на сваком дану користе само њихови свакодневно ажурирани
улази (одељак~\ref{subsec:ttm}).

Период 2014. године користи се као \textbf{година загревања}, у којој се прогнозе
производе, али се не вреднују, јер се на њима попуњава клизни прозор за оцену корекције повратне
трансформације и коефицијената рекалибрације. Евалуација обухвата период од 1. јануара
2015. до краја узорка, дакле 2015–2025.

\subsection{Корекција прогнозе}
\label{subsec:korekcija-prozor}

Непараметарска корекција $a$ из једначине~\eqref{eq:correction} и
коефицијенти Минцер–Зарновиц регресије (одељак~\ref{subsec:mz}) оцењују се на
\textbf{клизном прозору од последњих 252 дана стварних ван-узорачких грешака}, дакле
на прогнозама које је модел већ произвео за дане пре дана $t$, а чија је остварена
вредност у међувремену постала позната. Ове величине се стога никада не користе за
подешавање хиперпараметара нити за евалуацију, а по конструкцији не могу садржати
информацију из будућности.

Корекција се рачуна \textbf{обједињено преко симбола}, засебно за сваку комбинацију
модела, дизајна и хоризонта, будући да је реч о својству расподеле резидуала тог модела, а не
појединог симбола, па би оцена по симболу на кратком прозору била непотребно
нестабилна. Док прва остварена грешка још није доступна, односно на самом почетку периода
загревања, корекција остаје неутрална ($a = 0$). То одговара случају у коме се
повратна трансформација примењује без корекције.

\section{Евалуација}
\label{sec:evaluacija}

Метрике мере суштински различите ствари, па су разврстане у четири групе: мере
тачности прогнозе, мере статистичке значајности, економску евалуацију и
карактеризацију детекције скокова.

\subsection{Мере тачности прогнозе}

\paragraph{Фундаментални проблем.}
Стварна волатилност је немерљива, па се прогнозе пореде у односу на зашумљену заменску
меру.
Патон \parencite{patton2011} показао је да већина уобичајених функција губитка даје
\textbf{погрешно рангирање модела} када је заменска мера зашумљена, и извео потребне и довољне
услове које функција губитка мора испунити да би рангирање било робусно.

\paragraph{QLIKE (примарна метрика).}
Робусна је на шум у заменској мери:
\begin{equation}
\QLIKE_t = \frac{y_t}{\hat{y}_t} - \ln\!\left(\frac{y_t}{\hat{y}_t}\right) - 1
\label{eq:qlike}
\end{equation}

Губитак је једнак нули само када се прогноза тачно поклопи са оствареном вредношћу.
Функција је \textbf{асиметрична}, па се потцењивање волатилности кажњава јаче од
прецењивања, што је економски исправно, јер је потценити ризик скупље.

\paragraph{RMSE (секундарна).}
Такође робусна на шум у заменској мери, али симетрична:
\begin{equation}
\RMSE = \sqrt{\frac{1}{N}\sum_{t=1}^{N}\left(y_t - \hat{y}_t\right)^{2}}
\end{equation}

\paragraph{Средња апсолутна грешка (енгл.\ mean absolute error, MAE), додатна.}
Није робусна на шум у заменској мери:
\begin{equation}
\MAE = \frac{1}{N}\sum_{t=1}^{N}\left|y_t - \hat{y}_t\right|
\end{equation}

Пошто су QLIKE и RMSE минимизовани условном средином, а MAE условном медијаном, уз
MAE се обавезно рачуна и варијанта са корекцијом $a = 0$, да поређење по MAE не би
кажњавало модел зато што циља средину:
\begin{equation}
\MAE^{med} = \frac{1}{N}\sum_{t=1}^{N}\left|y_t - \exp\!\left(\widehat{\ln y}_t\right)\right|
\end{equation}

Ова варијанта рачуна се само за моделе оцењене у логаритамском простору, јер само они
имају прогнозу чију би Јенсенову корекцију требало поништити.

\paragraph{Неслагање рангирања.}
Редослед по QLIKE и по MAE сматра се стварно различитим само ако је
Диболд–Маријанов тест значајан на нивоу $0{,}05$ на обе метрике, и то по ћелији
коју чине хоризонт и дизајн.

\paragraph{Релативни QLIKE.}
Приказна мера која резултате чини одмах читљивим:
\begin{equation}
\text{relQLIKE} = \frac{\overline{\QLIKE}_{model}}{\overline{\QLIKE}_{log\_har}}
\label{eq:relqlike}
\end{equation}

Ово је \textbf{однос средина, никада средина односа}. Разлика није техничка ситница:
средина односа експлодира чим губитак референтног модела на једном једином дану приђе
нули, што се дешава, па један дан може преузети цео извештавани број. Референтни модел
је \texttt{log\_har}, а вредност $0{,}95$ значи 5\% боље од референтног.

\subsection{Мере статистичке значајности}
\label{subsec:nivo2-znacajnost}

\paragraph{Диболд–Маријанов тест.}
За два модела рачуна се дневна разлика губитака \parencite{dieboldMariano1995}:
\begin{equation}
d_t = L_t^{(A)} - L_t^{(B)}
\end{equation}

Тест статистика пореди просек те разлике са њеном дугорочном стандардном грешком:
\begin{equation}
\DM = \frac{\bar{d}}{\sqrt{\hat{\Omega}/N}}
\end{equation}

Пошто су циљне вредности на $h = 5$ и $h = 22$ преклапајући просеци, разлика губитака је
аутокорелисана, па се дугорочна варијанса $\hat{\Omega}$ оцењује Њуи–Вестовим
естиматором отпорним на хетероскедастичност и аутокорелацију (енгл.\ heteroskedasticity
and autocorrelation consistent, HAC) \parencite{neweyWest1987} са $h-1$ померајем, уз троугаоне тегове који
опадају са удаљеношћу помераја. На добијену статистику примењује се и Харви–Лејборн–
Њуболдова корекција за мале узорке \parencite{harvey1997}, која је скалира фактором
зависним од $N$ и $h$ и тиме уклања пристрасност HAC оцене при коначном узорку.

\paragraph{Свођење панела на дневну серију.}
Губици у овом раду чине панел, са по једном вредношћу за сваки пар (дан, симбол). Пре сваког
Диболд–Маријановог теста разлика губитака се најпре \textbf{усредњава преко симбола
унутар дана}, па се тек на тој дневној серији примењује HAC оцена са $h-1$ померајем.
Без тог корака померај $k$ не би значио један трговински дан. Пошто су у низу сложеном
по (дан, симбол) суседни елементи исти дан а различит симбол, сума у $\hat{\Omega}$
обухватила би попречни пресек унутар дана, док би стварна аутокорелација од
преклапајућих циљева на $h = 22$ лежала око $28 \times 21$ позиција даље и остала
невиђена. Истовремено би $N$ било надувано величином попречног пресека, па би
стандардна грешка била око $\sqrt{28}$ пута премала. Дневна серија уклања оба проблема
одједном, па је један померај један дан, $N$ је број дана, а попречна корелација међу
симболима апсорбована је у усредњавању уместо да се рачуна као независна информација.
Исти корак претходи и поступку из наредног пасуса, који као тест елиминације користи
управо Диболд–Маријанову статистику.

\paragraph{Контрола вишеструког тестирања.}
Са 36 дизајна и 16 модела, исцрпно поређење свих парова даје
хиљаде тестова, па би се при $\alpha = 0{,}05$ стотине „значајних“ резултата појавило
чисто случајно. Свака извештавана \textbf{породица} тестова стога пролази кроз
BH контролу стопе лажних открића, истом процедуром описаном у
одељку~\ref{subsec:a3}. Породица се мора навести експлицитно и бележи се уз резултат,
јер контрола стопе лажних открића нема смисла без изјаве над чим се контролише.

Унапред декларисана породица јесте \textbf{сваки модел наспрам референтног
модела log-HAR}, унутар исте ћелије (дизајн, хоризонт). Тиме се број поређења своди
на 15 по ћелији, уместо на све парове, а контраст одговара питању које рад
поставља.

\paragraph{Секвенцијална DM елиминација.}
Поступак је инспирисан радом Хансена, Лундеа и Насона \parencite{hansen2011}, који
полази од свих модела и итеративно одбацује најгори док се преостали не могу
статистички разликовати међусобно. Ниво значајности износи $\alpha = 0{,}10$.

Овде се, међутим, користи \textbf{упрошћена варијанта} и рад је тако и назива:
секвенцијална DM елиминација, а не скуп поверења модела. Изворни поступак гради
статистику ранга или полуквадратну статистику над \emph{свим} преживелима и њену
расподелу добија блок бутстрапом, чиме контролише грешку прве врсте целог низа
елиминација. У овом раду је тест елиминације једно Диболд–Маријаново поређење најгорег
са најбољим преживелим, без корекције за вишеструкост дуж низа. То је јефтинија
апроксимација која чува смисао поступка, али не носи његову формалну гаранцију;
ограничење је наведено у одељку~\ref{sec:ogranicenja}.

Треба нагласити да припадност преживелом скупу није рангирање тачкастих губитака:
елиминација зависи од односа разлике губитака према њеној стандардној грешци, па
стабилан модел који је благо лошији може остати, а модел са бољим просеком и
нестабилном серијом разлика бити одбачен. Унутар скупа нема уређења, а величина скупа
мери снагу узорка, а не квалитет модела.

\subsection{Декомпозиција варијансе}
\label{subsec:anova}

Ово је централни резултат рада и одељак који даје бројчан одговор на питање из наслова.

\textbf{Зависна променљива} јесте \textbf{дневни} QLIKE губитак, један ред по
комбинацији (модел, дизајн, хоризонт, симбол, дан). Коришћење средина ћелија не би
оставило никакав члан грешке, па оцена не би била могућа.

Пошто модели нису сви исте врсте у погледу зависности од дизајна мерења, спроводе се
\textbf{две одвојене декомпозиције}.

\paragraph{Декомпозиција (а) — модели осетљиви на дизајн.}
Обухвата моделе чији улази зависе од све три димензије мерења, дакле HAR-J, CHAR,
LightGBM, XGBoost, TTM-C и TTM-(C+J). Сваки од њих оцењен је на свих 36
дизајна, па су сви фактори мерења заступљени:
\begin{equation}
\begin{split}
L_{ijklmn} = {}& \mu + \tau_i^{model} + \gamma_j^{grid} + \delta_k^{est}
+ \psi_l^{test} + \rho_m^{regime} + \phi_n^{symbol} \\
& + (\gamma\delta)_{jk} + (\gamma\psi)_{jl} + (\tau\psi)_{il} + \varepsilon_{ijklmn}
\end{split}
\label{eq:anova}
\end{equation}

Ова декомпозиција даје први део одговора на хипотезу Х1, јер се пореди удео варијансе који
објашњавају три фактора мерења заједно са уделом који објашњава фактор модела.

\paragraph{Декомпозиција (б) — модели зависни само од мреже.}
Обухвата моделе чији улази не зависе од естиматора и теста (наивни, HAR, log-HAR,
SHAR, HARQ, HAR-IV, ARFIMA, MEM, \texttt{ttm\_rv}, \texttt{log\_har\_ttm}). Они су
оцењени једном по мрежи, па фактори естиматора и теста у њиховом случају уопште не
постоје:
\begin{equation}
L_{ijmn} = \mu + \tau_i^{model} + \gamma_j^{grid} + \rho_m^{regime}
+ \phi_n^{symbol} + \varepsilon_{ijmn}
\end{equation}

Ова декомпозиција даје други део одговора на хипотезу Х1, дакле да ли и сама фреквенција
узорковања, код модела који не виде ни естиматор ни тест, објашњава мерљив удео
варијансе губитка.

\paragraph{Зашто се две декомпозиције не обједињују.}
Модели из групе (б) немају никакву варијацију по естиматору и тесту. Њихово уношење у
једначину~\eqref{eq:anova} захтевало би да се исти ред понови кроз свих 12
комбинација естиматора и теста, чиме би у анализу ушле вештачки дуплиране, међусобно
зависне опсервације и исказала би се лажна прецизност. Ниједна од две декомпозиције
стога не садржи дуплиране редове, јер свака користи само моделе за које су сви њени
фактори стварно варирани. Јединствен збирни број није дефинисан и оба удела се
извештавају одвојено.

Четири својства ове поставке методолошки су суштинска.

\paragraph{1. Оцењивање по хоризонту, никада обједињено.}
Хоризонт стоји у кључу реда, али није фактор. Хоризонти имају различите скале циљне
величине, па би обједињавање бацило објашњиву варијансу у члан грешке и умањило сваки
парцијални ета-квадрат.

\paragraph{2. Симбол и режим као блокирајући фактори.}
Остављени у резидуалу, надували би члан грешке и смањили све извештаване уделе. Ознака
режима додељује се једном по пару (симбол, дан) и не зависи ни од дизајна ни од модела
(одељак~\ref{subsec:stratifikacija}), па се понаша као блок, једнако као симбол.

Пошто се граница режима рачуна обједињено преко свих симбола, ознака режима делом прати
ниво волатилности самог симбола, а делом стање тржишта тог дана, па на панелу овог рада
око трећине њене варијансе објашњава већ сам симбол. При суми квадрата типа III сваки
фактор добија само свој \emph{јединствени} допринос, па мали парцијални ета-квадрат
режима треба читати као допринос \emph{преко} онога што симбол већ објашњава, а не као
налаз да разлика између мирних и турбулентних дана не утиче на грешку прогнозе; та
разлика се посебно приказује у стратификацији (одељак~\ref{subsec:stratifikacija}).

\paragraph{3. Сума квадрата типа III са контрастима нултог збира (енгл. sum-to-zero).}
Дизајн је неуравнотежен, јер број важећих дана није исти за сваки симбол и дизајн,
расподела дана по режимима није равномерна, а модел садржи три структурне интеракције;
при таквом дизајну тип II и тип III не слажу се око главних ефеката. Тип III без
контраста нултог збира није добро дефинисан, па се контрастно кодирање
($\sum_i\tau_i = 0$ и тако редом) поставља експлицитно.

\paragraph{4. Интеракциони чланови су структурни, а не случајни.}
Микроструктурни шум подиже $\RV$ приближно линеарно са бројем интрадневних опсервација
(једначина~\ref{eq:noise-bias}), док естиматор отпоран на скокове тим шумом није
погођен на исти начин, па однос естиматора и $\RV$ заиста зависи од мреже. Изостављање
тих чланова пребацило би део ефекта у главне ефекте.

\textbf{Мера удела варијансе} јесте парцијални ета-квадрат:
\begin{equation}
\eta_p^{2} = \frac{SS_{effect}}{SS_{effect} + SS_{error}}
\label{eq:eta}
\end{equation}

\paragraph{Ограничење ове поставке.}
Различити дизајни мерења изведени су из \textbf{истих} сирових једноминутних свећица,
па су вредности $\RV$ на мрежама $1m$ и $5m$ за исти дан високо корелисане.
Претпоставка независности стандардног $F$-теста овде \textbf{не важи}. Последице су
следеће:

\begin{itemize}
\item Парцијални ета-квадрат извештава се као \textbf{дескриптиван} удео варијансе,
      што не захтева независност.
\item Сваки извештавани $F$ или $p$ мора се читати уз \textbf{двосмерно кластерисане}
      стандардне грешке \parencite{cgm2011}, кластерисане истовремено по дану и по
      симболу:
      \begin{equation}
      \hat{V}_{2way} = \hat{V}_{day} + \hat{V}_{symbol} - \hat{V}_{day \cap symbol}
      \end{equation}
\item Блок бутстрап преко дана, са дужином блока $\max(h,\,5)$ дана, предвиђен је као
      алтернативна инференцијална путања, али није спроведен због меморијских захтева
      (одељак~\ref{sec:rez-h1}); извештавана инференција почива на кластерисаним
      стандардним грешкама.
\end{itemize}

\paragraph{Мрежа мења и циљ, не само улаз.}
Губитак се на свакој мрежи рачуна у односу на реализовану варијансу исте мреже. Мреже
се зато не разликују само по улазима модела, него и по заменској мери у односу на коју се
прогноза оцењује: $\RV$ на петнаестоминутној мрежи има већу грешку узорковања, а на
једноминутној је више оптерећена микроструктурним шумом. QLIKE чува рангирање модела
при зашумљеној заменској мери \parencite{patton2011}, али не и ниво губитка, па удео фактора
мреже обухвата и различиту тежину самог циља, а не само утицај улаза на тачност. Удели
естиматора и теста за детекцију скокова тог проблема немају, јер се унутар исте мреже
циљ не мења.

\paragraph{Робусност на асиметрију је обавезна, а не опциона.}
Дневни QLIKE изразито је десно асиметричан, а узорак обухвата фебруар 2018. и период
пандемије, па мали број дана са екстремним губицима може одредити уделе варијансе док
излазна табела изгледа
сасвим уредно. Иста декомпозиција стога се изводи \textbf{четири начина, један поред
другог}, како је приказано у табели~\ref{tab:robusnost}.

\begin{table}[htbp]
\centering
\caption{Варијанте провере робусности декомпозиције варијансе}
\label{tab:robusnost}
\begin{tabularx}{\textwidth}{L{3.6cm} L{3.8cm} X}
\toprule
Варијанта & Зависна променљива & Сврха \\
\midrule
сирова                   & дневни QLIKE          & примарна \\
логаритамска             & $\ln(\QLIKE)$          & смањује утицај репова \\
рангови                  & ранг унутар дана      & потпуно непараметарска \\
без 1\%                   & без горњег 1\% дана   & одговор на питање да ли налаз важи
и без марта 2020. године \\
\bottomrule
\end{tabularx}
\end{table}

Стабилни удели кроз све четири варијанте представљају јаку тврдњу о робусности, док
се нестабилност наводи као налаз, а не прикрива.

Треба нагласити да тумачивост сирове варијанте зависи од репова расподеле. Ако
неколико дана носи већину укупног губитка, члан грешке постаје толико велик да ниједан
фактор не може добити осетан удео, па сирова варијанта тада не показује одсуство ефекта
него недостатак резолуције. Зато се логаритамска варијанта и варијанта са
ранговима овде унапред проглашавају варијантама на којима почива
провера робусности, при чему прва смањује утицај репова монотоном трансформацијом, а друга
одустаје од скале у целини. Сирова варијанта остаје примарна зависна променљива и
приказује се прва у свакој табели.

\subsection{Минцер–Зарновиц регресија}
\label{subsec:mz}

Регресија остварене вредности на прогнозу \parencite{mincerZarnowitz1969} открива
систематску пристрасност коју QLIKE не издваја посебно:
\begin{equation}
y_t = a + b\,\hat{y}_t + u_t
\end{equation}

Тестирају се хипотезе $a = 0$ и $b = 1$. Модел може имати низак губитак, а
доследно потцењивати ниво волатилности. Пошто су резидуали на дужим хоризонтима
аутокорелисани по конструкцији, користе се Њуи–Вестове HAC стандардне грешке са $h-1$
померајем.

Коефицијенти се оцењују \textbf{искључиво на клизном прозору ранијих ван-узорачких
грешака} (одељак~\ref{subsec:korekcija-prozor}), а рекалибрисани QLIKE извештава се
\textbf{упоредо} са сировим, никада уместо њега, јер би рекалибрација пре примарне
метрике представљала гледање унапред.

\subsection{Економска евалуација}
\label{subsec:var}

Прогноза волатилности користи се за израчунавање вредности при ризику (енгл.\ value at
risk, VaR):
\begin{equation}
\VaR_t(\alpha) = z_{1-\alpha}\sqrt{\hat{y}_t}
\end{equation}
где је $z_{1-\alpha}$ квантил стандардне нормалне расподеле ($2{,}33$ за ниво од 1\%).
Користи се квадратни корен прогнозиране варијансе, дакле стандардна девијација.
Економска евалуација спроводи се \textbf{само на хоризонту $h = 1$}, јер је само на
њему прогноза упоредива са појединачним дневним приносом; на $h = 5$ и $h = 22$ циљ је
просек преко више дана, па би поређење са једним приносом мешало временске скале.

\paragraph{Кључно методолошко усклађивање.}
Пошто $\RV$ не обухвата ноћни принос (одељак~\ref{subsec:rv}), $\sqrt{\hat{y}}$ је
стандардна девијација само интрадневног кретања, па се пробијања морају мерити у односу
на принос од отварања до затварања; серија од затварања до затварања систематски би
потценила ризик и учинила да тест одбацује хипотезу због неусклађености јединица.

Прогноза направљена на дан $t$ пореди се са приносом дана $t+1$ \textbf{истог
симбола}, где је $r_{d,t+1} = p_{t+1,n} - p_{t+1,0}$ логаритамски принос од цене прве
до цене последње свећице сесије тог дана. Индикатор пробијања дефинише се као
\begin{equation}
I_{t+1} = \mathbf{1}\left\{-r_{d,t+1} > \VaR_t(\alpha)\right\}
\end{equation}

\paragraph{Купијецов тест безусловне покривености.}
Проверава да ли број пробијања одговара декларисаном нивоу \parencite{kupiec1995}:
\begin{equation}
LR_{uc} = -2\ln\!\left[\frac{\alpha^{x}(1-\alpha)^{N-x}}
{\hat{\pi}^{x}(1-\hat{\pi})^{N-x}}\right]\ \sim\ \chi^{2}(1)
\end{equation}
где је $x$ број пробијања, а $\hat{\pi} = x/N$ њихов емпиријски удео. Случај без
иједног пробијања није изузетак него \textbf{одбацивање}, јер прогноза која никада није
пробијена систематски прецењује ризик, па се статистика рачуна уз конвенцију
$0\ln 0 = 0$ и даје велику вредност, а не нулу.

\paragraph{Кристоферсенов тест независности.}
Количником веродостојности над матрицом прелаза проверава да ли вероватноћа пробијања
зависи од тога да ли је претходни дан био пробијање \parencite{christoffersen1998};
модел са тачним бројем пробијања која се сва догоде исте недеље лош је модел ризика.

\paragraph{Динамички квантилни тест (енгл.\ dynamic quantile test, DQ).}
Строжа је верзија исте провере \parencite{engleManganelli2004}, у којој се центрирани индикатор
пробијања регресује се на сопствене претходне вредности и на ниво вредности при ризику,
па се тестира да ли су сви коефицијенти једнаки нули.

\paragraph{Лопезова функција губитка.}
Пробијања не само да броји, већ их кажњава сразмерно величини премашаја
\parencite{lopez1999}:
\begin{equation}
L^{Lopez} = \frac{1}{N}\sum_{t=1}^{N}
\left[1 + \left(\ell_{t} - \VaR_t\right)^{2}\right]I_t,
\qquad
\ell_{t} = -r_{d,t}
\label{eq:lopez}
\end{equation}
где $\ell_t$ означава губитак портфолија тог дана, дакле негативан принос, па позитивна
вредност значи пад цене. Ово разликује модел који греши мало од модела који греши
катастрофално, дакле разлику коју сам број пробијања не види.

Рангирање модела по Лопезовој функцији губитка пореди се са рангирањем по QLIKE
Кендаловим коефицијентом $\tau$, што представља непосредну проверу хипотезе Х3.

Тестови се рачунају на \textbf{обједињеној унији целог периода евалуације}
(приближно 2\,766 дана трговања \textbf{по симболу}), а не по краћим потпериодима. Период од
252 дана очекује свега око $2{,}5$ пробијања на нивоу од 1\%, што Купијецовом тесту
даје суштински никакву снагу.

\subsection{Карактеризација детекције скокова}

\paragraph{Методолошка напомена.}
Класификационе метрике (прецизност, одзив, $F_1$) овде \textbf{нису применљиве}, јер не
постоји основна истина, будући да скок није опсервабилан догађај, него оно што естиматор
прогласи скоком. Уместо тачности, мере се слагање и последице.

\paragraph{Жакаров индекс преклапања.}
Мери колико се две варијанте теста међусобно слажу:
\begin{equation}
J(A,B) = \frac{\left|A \cap B\right|}{\left|A \cup B\right|}
\end{equation}

Низак индекс значи да избор теста суштински мења улазне податке, што је сама премиса
рада.

\paragraph{Удео дана са скоком.}
Мери колико је метода конзервативна:
\begin{equation}
\text{удео} = \frac{\text{број дана са детектованим скоком}}{\text{укупан број дана}}
\end{equation}

Уобичајено навођен опсег од 5\% до 15\% \parencite{huangTauchen2005,abd2007} односи се
на \textbf{дневне} тестове примењене на \textbf{петоминутној} мрежи и ван те две
околности не важи, па одступање од њега није знак лоше подешеног прага. Ли–Микландов
тест тестира сваки интрадневни принос посебно и зато легитимно даје виши удео, а на
једноминутној мрежи и дневни тестови дају знатно веће уделе, јер микроструктурни шум подиже
$\RV$ приближно линеарно са бројем интрадневних опсервација
(једначина~\ref{eq:noise-bias}), док естиматор отпоран на скокове тим шумом није
погођен на исти начин, па разлика $\RV - \IVhat$ на којој тест почива садржи
систематску компоненту без везе са скоковима. Реч је о интеракцији димензија А1 и А3
коју рад уноси у декомпозицију варијансе као структурни члан
(одељак~\ref{subsec:anova}). Критеријум од 5\% до 15\% зато се примењује само на дневне
тестове на петоминутној мрежи.

\paragraph{Удео скоковите компоненте у укупној варијанси.}
Мери економску значајност:
\begin{equation}
\text{удео } J = \frac{\sum_t J_t}{\sum_t \RV_t}
\end{equation}

\paragraph{Плацебо тест.}
Индикатори скокова насумично се пермутују кроз време, чиме се задржава њихова
учесталост, али уништава веза са стварним данима. Пермутација се спроводи
\textbf{унутар симбола}, са 100 понављања, на дизајну 5m/BPV и хоризонту $h = 1$, а
стварна разлика у QLIKE-у пореди се са расподелом разлика добијених на пермутованим
подацима. Ако побољшање опстане и са пермутованим индикаторима, оно \textbf{није}
последица информације о скоковима, него последица структуре модела.

Важно ограничење јесте да је QLIKE недефинисан при $J_t = 0$ (дељење нулом и логаритам
нуле), па се никада не примењује на серију скоковите компоненте, већ искључиво на
прогнозу реализоване варијансе.

\subsection{Стратификација резултата}
\label{subsec:stratifikacija}

Просек кроз цео узорак може сакрити да модел побеђује у мирним, а губи у турбулентним
периодима, што је за управљање ризиком управо обрнуто од корисног. Све метрике рачунају
се одвојено по три димензије, приказане у табели~\ref{tab:stratifikacija}.

\begin{table}[htbp]
\centering
\caption{Димензије стратификације резултата}
\label{tab:stratifikacija}
\begin{tabular}{ll}
\toprule
Димензија & Подела \\
\midrule
Режим            & мирни наспрам турбулентних периода, по медијани $\RV_m$ \\
Присуство скока  & дани са скоком наспрам дана без скока \\
Хоризонт         & $h = 1$, $h = 5$, $h = 22$ \\
\bottomrule
\end{tabular}
\end{table}

Граница режима рачуна се \textbf{једном}, као медијана месечног агрегата $\RV_m$ на
петоминутној мрежи у првом периоду обуке, и трајно се памти; ознака се додељује по пару
(симбол, дан), па је иста за сваки дизајн и сваки модел. Када би се граница померала
како обрада напредује, ознака режима за дати дан постала би двосмислена, јер би зависила
од тога која је прогноза управо означавана.

\section{Преглед референтних радова}
\label{sec:odnos}

Литература која претходи овом раду развија се у два тока, који се ретко срећу.

Први ток пореди \textbf{моделе} при фиксираним мерама. Ту спадају радови који уводе
проширења хетерогеног ауторегресивног модела \parencite{corsi2009,abd2007,bpq2016} и
радови који примењују методе машинског и дубоког учења
\parencite{christensen2023,leushuisPetkov2026}. Они проширују скуп архитектура, али
улазне реализоване мере узимају као дате, па сви модели у поређењу виде исти, унапред
изабран скуп мера.

Други ток пореди \textbf{мере} при фиксираном моделу. Ту спадају радови који уводе
нове естиматоре и тестове за детекцију скокова
\parencite{bns2004,andersen2012,corsiPirinoReno2010}, као и систематско поређење
реализованих мера као заменских мера \parencite{lps2015}. Они пажљиво испитују својства саме
мере, али последице избора мере обично не прате кроз више модела и хоризоната.

Изузетак у односу на први ток представља HARQ модел из одељка~\ref{subsec:harq}, који
узима у обзир да мерење садржи грешку и да та грешка варира кроз време, али је моделује
\textbf{унутар} једног модела, уместо да је третира као димензију која се може
варирати.

Овај рад спаја оба тока на истом панелу, тако што се мере и модели варирају истовремено, у
потпуном факторском плану, а декомпозиција варијансе из одељка~\ref{subsec:anova}
раздваја уделе варијансе губитка који припадају једном и другом извору. Тиме се питање
колико сваки од њих доприноси тачности исказује као број, а не као оцена.

Апарат описан у овом поглављу чини оквир унутар кога се централно питање рада може
бројчано проверити, и то архитектура модела као прва оса, дизајн мерења као друга, изведене
променљиве које их повезују и декомпозиција варијансе као мерило. Начин на који
је тај оквир спроведен описан је у поглављу~\ref{ch:metodologija}.

\chapter{Методологија}
\label{ch:metodologija}

Ово поглавље описује поступак којим је спроведен експеримент. Најпре је дат преглед
целокупног истраживачког тока, затим извор и припрема података, потом дизајн
експеримента са свим фиксираним параметрима и, на крају, мере, правила искључивања,
план статистичке анализе и ограничења методологије. Све вредности наведене у овом
поглављу идентичне су у свим условима експеримента, осим оних које саме представљају
експерименталне факторе.

\section{Преглед истраживачког поступка}

Истраживање је спроведено као факторски експеримент у коме се укрштају две
експерименталне осе описане у поглављу~\ref{ch:teorija}: дизајн мерења и архитектура
модела. Циљ поступка јесте да се разлике у тачности предвиђања припишу појединачним
изворима, па су сви кораци подређени захтеву да свака комбинација буде обрађена под
идентичним условима.

Поступак се одвија у четири узастопне фазе:

\begin{enumerate}
\item \textbf{Прикупљање и чишћење података.} Интрадневне свећице на једноминутној
      резолуцији преузимају се од добављача London Strategic Edge, филтрирају према стварном
      календару трговања и проверавају на потпуност.
\item \textbf{Конструкција панела мерења.} За сваку од тридесет шест комбинација
      дизајна мерења рачунају се реализоване мере, спроводи се условна декомпозиција
      на континуирану и скоковиту компоненту и граде се HAR агрегати.
\item \textbf{Моделовање.} Сваки модел се оцењује и предвиђа на проширујућем прозору,
      уз поновно оцењивање сваких 21 дан трговања, за сваки хоризонт и сваки дизајн од
      кога зависи.
\item \textbf{Евалуација.} Прогнозе се оцењују функцијама губитка, пореде тестовима
      значајности, подвргавају економској провери и разлажу декомпозицијом варијансе.
\end{enumerate}

Дијаграм целокупног тока приказан је на слици~\ref{fig:pipeline}.

\begin{figure}[htbp]
\centering
\begin{tikzpicture}[node distance=9mm]

\node[faza] (izvor) {\textbf{Извор}\\[2pt] једноминутне свећице, 28 симбола};
\node[faza, below=of izvor] (ciscenje) {\textbf{Чишћење}\\[2pt] календар трговања, покривеност, корпоративне акције};
\node[faza, below=of ciscenje] (panel) {\textbf{Панел мерења}\\[2pt] реализоване мере и условна декомпозиција};
\node[faza, below=of panel] (modeli) {\textbf{Моделовање}\\[2pt] проширујући прозор, оцењивање на 21 дан};
\node[faza, below=of modeli] (eval) {\textbf{Евалуација}\\[2pt] четири нивоа метрика};

\draw[tok] (izvor) -- (ciscenje);
\draw[tok] (ciscenje) -- (panel);
\draw[tok] (panel) -- (modeli);
\draw[tok] (modeli) -- (eval);

% десна колона: шта се у којој фази варира, односно добија
\node[detalj, right=13mm of ciscenje] (d1)
  {исти скуп дана на све три мреже};
\node[detalj, right=13mm of panel] (d2)
  {А1 $\times$ А2 $\times$ А3 \\ \textbf{36 дизајна}};
\node[detalj, right=13mm of modeli] (d3)
  {\textbf{16 модела} \\ $h \in \{1, 5, 22\}$};
\node[detalj, right=13mm of eval] (d4)
  {тачност, значајност, економска провера, скокови};

\draw[tok, dashed, gray] (ciscenje) -- (d1);
\draw[tok, dashed, gray] (panel) -- (d2);
\draw[tok, dashed, gray] (modeli) -- (d3);
\draw[tok, dashed, gray] (eval) -- (d4);

% лева колона: период
\node[medjukorak, left=13mm of panel] (p1) {2010–2025.\\ изградња панела};
\node[medjukorak, left=13mm of modeli] (p2) {2014.\\ загревање};
\node[medjukorak, left=13mm of eval] (p3) {2015–2025.\\ евалуација};

\draw[tok, dashed, gray] (p1) -- (panel);
\draw[tok, dashed, gray] (p2) -- (modeli);
\draw[tok, dashed, gray] (p3) -- (eval);

\end{tikzpicture}
\caption{Ток обраде података и оцењивања модела}
\label{fig:pipeline}
\end{figure}

\section{Подаци}

\subsection{Извор података}

Полазна напомена тиче се доступности извора. Збирка \textit{Oxford-Man Realized
Library}, која је дуго била стандардни извор у овој области, више није доступна и
институција је саопштила да не планира њену замену \parencite{oxfordMan2023}. Пошто
готово сви радови на које се овај рад позива користе управо њу, њихова поставка се не
може непосредно поновити, па је извор података морао бити изабран самостално.

Интрадневне свећице на једноминутној резолуцији преузимају се масовним извозом од
добављача \textit{London Strategic Edge}; уз сваку преузету датотеку бележи се порекло,
тако да је преузимање поновљиво.

Индекс имплицитне волатилности преузима се из базе \textit{FRED} Федералних резерви
Сент Луиса, серија \texttt{VIXCLS}.

\subsection{Узорак}
\label{subsec:uzorak}

Узорак обухвата \textbf{двадесет осам акција} којима се тргује на америчким берзама,
у периоду од 1. јануара 2010. до 31. децембра 2025. године. Симболи су изабрани према
дужини расположиве историје и обиму трговања, уз захтев да свака серија има најмање
шеснаест година једноминутних података.

Првобитни избор обухватао је тридесет симбола, од којих су два искључена: једна током
контроле квалитета, пре моделовања, а друга након пробног покретања.

Акција компаније ArcelorMittal (MT) искључена је током контроле квалитета, зато што се у двогодишњем
периоду пре емисије права пречег уписа из марта 2016. године цене унутар истог дана
смењују између прилагођеног и неприлагођеног нивоа. Последица су привидни
петоминутни приноси од око 23\%, што одговара управо фактору прилагођавања те
емисије, на 496 дана. Такви подаци произвели би лажне скокове и вишеструко надуване
вредности реализоване варијансе, па се серија не може користити ни делимично без
произвољног избора периода.

Акција компаније NVIDIA (NVDA) искључена је касније, \textbf{након пробног покретања},
и то не због сумње у податке него због онога што је покретање показало: та једна серија
носила је 99,77\% укупне масе QLIKE губитка у евалуацији, па је просечан губитак
наивног модела са њом износио 14,161, а без ње 2,32. Провером узрока утврђено је да је
реч о \textbf{дискретности цене}. Цене у извору уназад су прилагођене за две поделе акција,
у односу 4:1 из 2021. и 10:1 из 2024. године, па се у периоду 2014–2016. серија
креће око $0{,}46$ америчких долара. Пошто се цене на берзи бележе са тачношћу од
једног цента, најмањи могући корак цене у том периоду износи приближно $2{,}2\%$
вредности акције. Реализована варијанса тада не мери волатилност него заокруживање:
у панелу постоји 65 дана са $\RV_d = 0$, сви на тој серији, а однос $\RV(1m)/\RV(5m)$
за те године износи око $2{,}3$, наспрам $1{,}14$ колико износи на целом панелу.
Пошто би губитак те једне серије сам одредио сваки извештавани просек, а мерио би
заокруживање а не волатилност, серија је искључена у целини. Одлука је донета након
пробног, а пре коначног покретања, и наведена је овде управо зато што није била
утврђена унапред.

Панел уместо једне серије изабран је из три разлога. Прво, једна серија од шеснаест
година даје око четири хиљаде дневних опсервација, што је недовољно за стабилно
оцењивање модела са већим бројем параметара, какви су модели дубоког учења, који су у
литератури тестирани на скуповима са стотинама хиљада тачака, док панел даје ред
величине више. Друго, налаз потврђен на
двадесет осам симбола робуснији је од налаза на једној серији. Треће, панел омогућава
проверу да ли закључци зависе од ликвидности и сектора.

Четири од двадесет осам симбола односе се на компаније које нису основане у Сједињеним
Државама, али се њиховим акцијама тргује на америчким берзама. То је наведено
ограничење, а не замена за неамеричко тржиште: сви симболи деле исти календар
трговања, па рад не проверава робусност налаза на тржишту са другачијим радним
временом.

\subsection{Чишћење и контрола квалитета}
\label{subsec:ciscenje}

Свакој свећици приступа се кроз низ провера изведених из \textbf{стварног календара
трговања}, а не из фиксираног временског прозора. Фиксиран прозор би на данима са
ранијим затварањем задржао велики број танких свећица након затварања тржишта, чији
је медијана обима реда величине четрдесет пута мања од обима у правој сесији.

Примењују се следећи кораци:

\begin{enumerate}
\item \textbf{Филтрирање према сесији.} Задржавају се само свећице унутар календарског
      прозора сесије за тај конкретан дан. Дани који нису права берзанска сесија
      одбацују се у целини, јер свећице од добављача повремено упадају на нерадне дане.
\item \textbf{Одбацивање артефаката.} Уклањају се свећице које истовремено имају нулти
      обим и нулту промену цене. Свећица са стварним обимом али без промене цене
      задржава се, јер представља легитиман миран тренутак.
\item \textbf{Уклањање дупликата} и сортирање по временској ознаци.
\item \textbf{Искључивање скраћених дана.} Дани са ранијим затварањем (210 минута
      трговања) искључују се у целини, са \textbf{свих} мрежа.
\item \textbf{Провера покривености.} Дан се задржава само ако број једноминутних
      свећица достиже најмање 80\% очекиваног броја за пун дан (390) и ако свака
      мрежа на том дану даје најмање двадесет приноса.
\end{enumerate}

Оба правила служе истом циљу: све три мреже морају да деле \textbf{исти скуп дана},
а праг се примењује једном, на нивоу целе мреже, а не засебно за сваки естиматор. Када
би се скраћени дан задржао на мрежама $1m$ и $5m$ а одбацио на мрежи $15m$, или када би
се дан задржао за један естиматор а одбацио за други, фактор А1 односно А2 мењао би и
скуп дана над којим се рачуна губитак, па удели варијансе у декомпозицији не би били
упоредиви.

Поред тога, за сваки симбол спроводи се \textbf{тврда провера прилагођености
корпоративним акцијама}: на дан сваке поделе акција пореди се стварни дневни принос са
падом који би се видео да цене нису прилагођене, а серија се проглашава неприлагођеном
ако је принос ближи том паду него нули. Тада се изградња панела за тај симбол
зауставља, јер би неприлагођена серија произвела лажне скокове. Управо је та провера
издвојила серију MT.

\subsection{Конструкција панела}

За сваки симбол и сваку од тридесет шест комбинација дизајна мерења рачунају се
реализоване мере и условна декомпозиција, а HAR агрегати граде се клизањем преко
\textbf{доступних дана трговања}, никада календарских, при чему се непотпуни прозори
остављају као недостајући.

Пре уписивања, над састављеним панелом спроводе се две тврде провере исправности:
идентитет семиваријанси из једначине~\eqref{eq:semivar-identity} и поштовање
апсолутног прага броја интрадневних опсервација. За разлику од бројаних дијагностика,
ове провере заустављају процес ако не прођу, јер њихов пад указује на грешку у рачуну,
а не на својство података.

\section{Дизајн експеримента}

\subsection{Фактори и нивои}

Експеримент укршта две експерименталне осе, чији су нивои дати у
табели~\ref{tab:nivoi}.

\begin{table}[htbp]
\centering
\caption{Експерименталне осе, нивои и зависност модела од дизајна мерења}
\label{tab:nivoi}
\begin{tabularx}{\textwidth}{L{2.6cm} X c}
\toprule
Оса или група & Нивои, односно модели & Број \\
\midrule
А1 — фреквенција узорковања & 1m, 5m, 15m & 3 \\
А2 — естиматор отпоран на скокове & BPV, MedRV, TrBPV & 3 \\
А3 — тест за детекцију скокова & наивни, BNS, LM, LM-FDR & 4 \\
\addlinespace
Модели зависни само од А1 &
наивни, HAR, log-HAR, SHAR, HARQ, HAR-IV, ARFIMA, MEM, ttm-rv, log-HAR-TTM &
10 \\
\addlinespace
Модели зависни од А1, А2 и А3 &
HAR-J, CHAR, LightGBM, XGBoost, ttm-c, ttm-c+j &
6 \\
\bottomrule
\end{tabularx}
\end{table}

Комбиновањем трију димензија дизајна мерења добија се тридесет шест дизајна, а укупан
број модела износи шеснаест. Свака комбинација оцењује се на три хоризонта
($h = 1$, $h = 5$ и $h = 22$ дана).

\subsection{Свођење рачунског обима}
\label{subsec:svodjenje}

Оцењивање сваког модела на сваком дизајну било би непотребно скупо, јер већина модела
не зависи од свих трију димензија дизајна мерења. Модел који као улаз користи само
реализовану варијансу даје идентичне прогнозе под свих дванаест комбинација естиматора
и теста на истој мрежи, па би његово поновно оцењивање било само трошење времена.

Сваки модел зато декларише \textbf{од којих димензија његови улази стварно зависе}, а
систем га оцењује једном по свакој различитој комбинацији само тих димензија: десет
модела из табеле~\ref{tab:nivoi} оцењује се на три дизајна, а преосталих шест на свих
тридесет шест.

Број оцењених комбинација стога износи $(10 \times 3 + 6 \times 36)$ дизајна
$\times$ 3 хоризонта, поново оцењивано сваких 21 дан трговања.

Ова одлука није само рачунска. Понављање идентичних редова кроз дизајне од којих модел
не зависи унело би вештачки дуплиране опсервације у декомпозицију варијансе и исказало
би лажну прецизност; управо зато се декомпозиције (а) и (б) из
одељка~\ref{subsec:anova} спроводе одвојено.

\subsection{Фиксирани параметри}

Све вредности које нису предмет варирања фиксиране су пре покретања експеримента и
нису подешаване према резултатима. Све вредности наведене су у табели~\ref{tab:fiksirani}, у Прилогу А.

\subsection{Провера контаминације предтренирања}
\label{subsec:kontaminacija}

Пре него што се произведе иједна прогноза претходно обученог модела, спроводи се провера у којој
се објављени датум одсецања података за предтренирање пореди са периодом евалуације.
Резултат се уписује као трајан артефакт \textbf{у оба случаја}, и када преклапања нема
и када га има.

Када се преклапање утврди, оно се квантификује, а модел се задржава у
анализи уз изричиту ознаку да му период евалуације делимично претходи датуму
одсецања. Модел се никада не избацује тихо нити се тихо третира као чист, јер би оба
поступка учинила резултат непроверљивим.

У спроведеном експерименту провера је утврдила преклапање: од 2\,766 трговинских дана
у периоду евалуације, њих 2\,264 (81,9\%, дакле период 2015–2023) претходи објављеном
датуму одсецања од 1. јануара 2024. године, док је период 2024–2025. у целини чист.
Реч је о преклапању периода, а не о доказаном цурењу података, јер није познато да ли
су у скуп за предтренирање ушле управо серије ових симбола.

\section{Мере}

Мере су унапред подељене у три групе, како би се избегло накнадно тражење статистички
значајних разлика међу великим бројем излазних величина. \textbf{Примарне} су QLIKE и
релативни QLIKE, удели објашњене варијансе из обе декомпозиције и припадност скупу
поверења модела. \textbf{Секундарне} су RMSE и MAE са варијантом за условну медијану,
коефицијенти Минцер–Зарновиц регресије и рекалибрисани QLIKE, резултати тестова
вредности при ризику и Кендалов коефицијент слагања рангирања. \textbf{Дијагностичке}
мере не служе за проверу хипотеза него за контролу исправности: удео дана са скоком и
индекс преклапања по дизајну, удео скоковите компоненте, број искључених дана по
разлогу и број опсервација по ћелији. Све су дефинисане у одељку~\ref{sec:evaluacija}.

\section{Правила искључивања}

Правила наведена испод утврђена су, по правилу, пре коначног покретања и пре увида у
његове резултате; одступање од тог редоследа представља искључење серије NVDA,
објашњено у одељку~\ref{subsec:uzorak}, до кога је довело пробно покретање:

\begin{itemize}
\item дани у периоду загревања клизних прозора не улазе у статистику о скоковима, јер
      за њих локална оцена волатилности није дефинисана;
\item дани за које циљна променљива није дефинисана, дакле последњих $h$ дана сваке
      серије, искључују се из евалуације на том хоризонту;
\item симбол чија серија у тренутку оцењивања има премало опсервација изоставља се из
      оцењивања модела који се оцењују засебно по симболу, уз бележење; праг износи 60
      опсервација за ARFIMA, код кога оцена дугог памћења и избор редова захтевају дужу
      серију, и 30 опсервација за MEM;
\item редови који се не могу оценити не улазе у метрике, али се броје и приказују по
      моделу и разлогу. То су дани на којима је остварена циљна вредност једнака нули,
      где QLIKE није дефинисан, и дани на којима прогноза недостаје или није позитивна,
      укључујући и случајеве у којима нумеричка заштита ARFIMA модела не да употребљиву
      вредност;
\item резултат се не искључује због тога што одступа од очекиваног. Дани са
      изразито великим губитком задржавају се у примарној анализи, а њихов утицај се
      проверава засебном варијантом декомпозиције из које је изузет горњи проценат дана
      по губитку.
\end{itemize}

\section{Статистичка анализа}
\label{sec:analiza}

Начин провере сваке хипотезе из одељка~\ref{sec:hipoteze} дат је у
табели~\ref{tab:plan}. Све вредности значајности прате правила из
одељка~\ref{sec:evaluacija}: свака извештавана породица поређења подвргава се контроли
стопе лажних открића и наводи се експлицитно, уз сваку значајност наводи се и величина
ефекта, а закључак се никада не изводи само из вредности значајности.

\begin{table}[htbp]
\centering
\caption{План провере хипотеза}
\label{tab:plan}
\begin{tabularx}{\textwidth}{L{1.1cm} L{4.3cm} X L{1.9cm}}
\toprule
Хип. & Поступак & Скуп модела и услови & Резултати \\
\midrule
Х1 &
декомпозиција варијансе дневних губитака, одвојено по хоризонту, симбол и режим као
блокови &
(а) модели осетљиви на дизајн, свих 36 дизајна; (б) модели зависни само од мреже, три
дизајна; удели се извештавају одвојено &
одељак~\ref{sec:rez-h1} \\
\addlinespace
Х2 &
Диболд–Маријанов тест над дневним губицима, уз контролу стопе лажних открића &
парови HAR-J и CHAR наспрам log-HAR, унутар истог дизајна и хоризонта; допунски
поређење величине ефекта по хоризонтима и плацебо тест са 100 пермутација унутар
симбола (5m/BPV, $h = 1$) &
одељак~\ref{sec:rez-h2} \\
\addlinespace
Х3 &
Кендалов коефицијент $\tau$ између рангирања по QLIKE и по Лопезовој функцији губитка &
сви модели, само хоризонт $h = 1$, тестови вредности при ризику на целом периоду
евалуације &
одељак~\ref{sec:rez-h3} \\
\bottomrule
\end{tabularx}
\end{table}

Декомпозиција варијансе спроводи се у четири варијанте наведене у
табели~\ref{tab:robusnost}, једна поред друге; стабилност удела кроз све четири
варијанте основ је за тврдњу о робусности, док се нестабилност наводи као налаз.

\section{Рачунско окружење и поновљивост}

Цео поступак имплементиран је у програмском језику Python, уз управљање зависностима
које фиксира тачне верзије свих коришћених библиотека: NumPy, pandas, statsmodels,
scikit-learn, LightGBM, XGBoost, transformers и exchange\_calendars. Изворни код
доступан је на захтев. Естиматори, тестови и
модели разрешавају се преко назива из централног регистра, па ниједна варијанта не
може добити повлашћен третман у коду, а аутоматизовани тестови покривају распоред
оцењивања, повратну трансформацију и одсуство преласка између симбола при грађењу
циљне променљиве. Ниједан модел није стохастички, па поновљено покретање над истим
подацима даје идентичне прогнозе; подаци се не чувају уз код, јер су поново изводиви
из извора.

\section{Ограничења методологије}
\label{sec:ogranicenja}

Следећа ограничења проистичу из самог дизајна истраживања и наводе се унапред.

\begin{itemize}
\item \textbf{Зависност између дизајна.} Свих тридесет шест дизајна изведено је из
      истих сирових једноминутних свећица, па су вредности реализоване варијансе за
      исти дан високо корелисане. Претпоставка независности класичног теста стога не
      важи, због чега се удели варијансе исказују као дескриптивни, а свака значајност
      чита уз кластерисане стандардне грешке.
\item \textbf{Пристрасност преживљавања.} Симболи су изабрани по захтеву да имају пуну
      шеснаестогодишњу историју једноминутних података, што унапред искључује компаније
      које су у међувремену престале да се котирају, припојене су или су банкротирале.
      Узорак стога представља преживеле, а не случајан пресек тржишта.
\item \textbf{Један добављач података.} Све интрадневне свећице потичу од једног
      извора, па се грешке својствене том извору не могу открити поређењем са другим
      добављачем.
\item \textbf{Дискретност цене.} Реализоване мере претпостављају да је цена непрекидна
      величина, док се на берзи бележи са тачношћу од једног цента. При ниским ценама
      акције тај корак чини осетан део вредности, па реализована варијанса делом мери
      заокруживање, а не волатилност. У овом раду то је водило искључењу серије NVDA,
      али исти механизам у мањој мери погађа сваку серију са ниском ценом, због чега се
      налази не могу без провере пренети на акције које се тргују по ниским ценама.
\item \textbf{Изостављени естиматори отпорни на шум.} Димензија А2 обухвата естиматоре
      отпорне на скокове, али не и оне отпорне на микроструктурни шум, попут
      реализованог језгра или поступка претходног усредњавања. Питање колики би удео
      варијансе припао таквим естиматорима остаје отворено.
\item \textbf{Једно тржиште.} Сви симболи деле исти календар трговања, па рад не
      проверава да ли закључци важе на тржишту са другачијим радним временом.
\item \textbf{Искључен ноћни принос.} Реализована волатилност обухвата само
      интрадневно кретање, па се и економска евалуација односи искључиво на интрадневни
      ризик.
\item \textbf{Изостављени модели.} Скуп модела не обухвата све архитектуре из
      литературе; избор је ограничен на оне које су или стандард у области или
      средиште актуелне расправе. Трансформерска архитектура PatchTST није укључена
      зато што је панел премали да би се трансформер обучио од нуле, а трошак обучавања
      био би несразмеран доприносу; правац дубоког учења у раду представља претходно
      обучени модел опште намене, а обучавање трансформера на већем панелу остаје
      предмет даљег истраживања.
\item \textbf{Упрошћен скуп поверења модела.} Поступак елиминације описан у
      одељку~\ref{subsec:nivo2-znacajnost} користи једно Диболд–Маријаново поређење
      најгорег са најбољим преживелим моделом, без корекције за вишеструкост дуж низа
      елиминација. Изворни поступак Хансена и сарадника уместо тога гради статистику
      над свим преживелима и њену расподелу добија блок бутстрапом. Пријављена
      припадност стога јесте у духу скупа поверења модела, али не носи његову формалну
      гаранцију, па се разлике у величини скупа између ћелија читају уздржано.
\item \textbf{Фиксирани хиперпараметри.} Модели засновани на стаблима нису подешавани,
      па њихови резултати представљају доњу границу могућег учинка те архитектуре. Та
      одлука је намерна, јер би подешавање помешало уложен труд у претрагу са ефектом
      саме архитектуре.
\item \textbf{Преклапање периода предтренирања.} Период евалуације 2015–2023. претходи
      објављеном датуму одсецања података за предтренирање коришћеног модела опште
      намене, док је период 2024–2025. чист. Реч је о преклапању периода, а не о
      доказаном цурењу података, али се резултати тог модела на ранијем делу узорка не
      могу тумачити као потпуно ван-узорачки.
\end{itemize}

\chapter{Резултати}
\label{ch:rezultati}

Резултати експеримента биће приказани по завршетку пуног покретања.

\section{Узорак и карактеризација скокова}
\label{sec:rez-uzorak}

\section{Тачност предвиђања}
\label{sec:rez-tacnost}

\section{Хипотеза Х1 — декомпозиција варијансе}
\label{sec:rez-h1}

\section{Хипотеза Х2 — информација о скоковима}
\label{sec:rez-h2}

\section{Хипотеза Х3 — економска евалуација}
\label{sec:rez-h3}

\section{Стратификација резултата}
\label{sec:rez-stratifikacija}

\section{Дискусија}
\label{sec:rez-diskusija}

\chapter{Закључак}
\label{ch:zakljucak}

\section{Главни налази}

\section{Практични значај}

\section{Правци даљег истраживања}

Природни наставак рада чини неколико праваца, од којих сваки одговара по једном
ограничењу наведеном у одељку~\ref{sec:ogranicenja}:

\begin{itemize}
\item проширење димензије А2 естиматорима отпорним на микроструктурни шум, попут
      реализованог језгра и поступка претходног усредњавања;
\item провера налаза на тржиштима са другачијим радним временом и на другим класама
      имовине;
\item укључивање ноћног приноса кроз мере које обухватају целу календарску промену
      цене, чиме би се економска евалуација проширила на укупан дневни ризик;
\item укључивање трансформерских архитектура обучених од нуле, на панелу довољно
      великом да такво обучавање буде оправдано;
\item провера модела опште намене на издању чији датум одсецања података у целини
      претходи периоду евалуације.
\end{itemize}


% ── литература (поглавље се НЕ нумерише, по шаблону) ────────
\printbibliography[heading=bibintoc,title={ЛИТЕРАТУРА}]

% ── прилози ─────────────────────────────────────────────────
\chapter*{ПРИЛОГ А — ФИКСИРАНИ ПАРАМЕТРИ}
\addcontentsline{toc}{chapter}{Прилог А — фиксирани параметри}
\label{app:parametri}

Табела обухвата све параметре који су фиксирани пре покретања експеримента и који се
не мењају ни у једном услову, осим оних који сами представљају експерименталне факторе.

\begin{center}
\begin{longtable}{L{4.4cm} L{3.4cm} L{5.2cm}}
\caption{Фиксирани параметри експеримента}
\label{tab:fiksirani}\\
\toprule
Параметар & Вредност & Разлог фиксирања \\
\midrule
\endfirsthead
\toprule
Параметар & Вредност & Разлог фиксирања \\
\midrule
\endhead
\midrule
\multicolumn{3}{r}{\footnotesize наставак на следећој страни}\\
\endfoot
\bottomrule
\endlastfoot
Период узорка & 2010–2025 & расположивост једноминутне историје \\
Период евалуације & 2015–2025 & 2014. је година загревања клизних прозора \\
Праг за TrBPV & 3,0 дневне робусне девијације & увођење као четврте димензије дало би неуравнотежен план \\
Праг наивног теста & 3 стандардне девијације, двострано & доња граница квалитета, није озбиљна метода \\
Ниво значајности BNS теста & $\alpha = 0{,}001$, једнострано & строг праг смањује број лажних скокова у дугом узорку \\
Ниво Гумбелове критичне вредности (LM) & 0,999 & контрола грешке унутар дана преко свих интрадневних тестова \\
Стопа лажних открића (LM-FDR) & $q = 0{,}05$ & контролише удео лажних међу проглашеним данима; при овим нивоима блаже од Гумбеловог прага \\
Локални прозор LM теста & 4 сесије на свакој мрежи & да се са мрежом мења само фреквенција, а не и дужина прозора \\
Клизни прозор FDR корекције & 250 дана & приближно једна година трговања; спречава гледање унапред \\
Минимална покривеност дана & 80\% очекиваног броја једноминутних свећица & филтрирање непотпуних сесија, иста одлука на свакој мрежи \\
Апсолутни праг опсервација & 20 приноса на датој мрежи & обезбеђује стабилну оцену реализоване квартичности \\
Поновно оцењивање модела & сваких 21 дан трговања & приближно месечно; улази се ажурирају свакодневно \\
Клизни прозор за корекцију и рекалибрацију & 252 дана ван-узорачких грешака & приближно једна година; искључиво прошле грешке \\
Хиперпараметри стабала & 200 стабала, дубина 4, 15 листова, стопа 0,05 & подешавање би мешало труд у претрази са ефектом архитектуре \\
Контролна тачка претходно обученог модела & \texttt{ibm-granite/\allowbreak granite-\allowbreak timeseries-\allowbreak ttm-r2}, издање \texttt{512-192-r2} & различита издања имају различите скупове предтренирања \\
Тежине хибридног модела & 0,5 и 0,5 & оптимизација би одразила претрагу, не архитектуру \\
Ниво скупа поверења модела & $\alpha = 0{,}10$ & уобичајен ниво у литератури о скупу поверења \\
Ниво Диболд–Маријановог теста & $\alpha = 0{,}05$, уз контролу стопе лажних открића & велики број поређења захтева контролу породице \\
Граница режима & медијана $\RV_m$ у првом периоду обуке, мрежа 5m & фиксна граница чини ознаку режима једнозначном \\
Ширина појаса локалног Витловог естиматора & $m = \lfloor n^{0{,}65}\rfloor$ & уобичајен избор облика $n^{\alpha}$ \\
Број пермутација у плацебо тесту & 100, унутар симбола & расподела разлика под нултом хипотезом \\
Ниво вредности при ризику & 1\% & стандард у регулаторној пракси \\
\end{longtable}
\end{center}


\end{document}
